% Research fragment prepared by /root/rectangular_bridge in run CUE_GCT-20261121.
% Requires amsmath, amssymb, amsthm; no nonstandard macros are used.
% Citations use bibliographic keys listed in sources.md.
% These are checked deductions in this run, not assertions of literature priority.

\section{Rectangular characters, explicit positive witnesses, and orbit ranks}

Throughout this section the ground field is $\mathbb C$, compact-group Haar
measure has mass one, and $D=nd$. We write $r=(d^n)$ and
$D_\alpha(b)=s_\alpha(1^b)=\dim S_\alpha(\mathbb C^b)$.
Duals are displayed when coordinate rings are involved. Passing from the
polynomial representation to the coordinate representation changes both sides
by the same duality and leaves every multiplicity below unchanged.

\subsection{R-01: the exact rectangular restriction transform}

\paragraph{Statement R-01.}
Let $a\geq\ell(\lambda)$ and let $X=(x_1,\ldots,x_a)$,
$Y=(y_1,\ldots,y_n)$. Denote the product alphabet by
$XY=(x_i y_j)_{1\leq i\leq a,1\leq j\leq n}$. Then
\begin{equation}\label{R01:branching}
 s_r(XY)=\sum_{\mu,\nu\vdash D}
 g(r,\mu,\nu)s_\mu(X)s_\nu(Y).
\end{equation}
In particular, by symmetry of the three Kronecker labels,
\begin{equation}\label{R01:haar}
 g(\lambda,r,r)=\int_{U(a)\times U(n)}
 s_r(A\otimes B)\overline{s_\lambda(A)}\det(B)^{-d}\,dA\,dB.
\end{equation}
Every coefficient needed on the rectangular-Kronecker side of determinant
GCT is therefore an explicit tensor-subgroup Fourier coefficient of the
\emph{same rectangular character} that occurs in the finite CUE moment.
The equality is not an identification of that coefficient with the scalar
moment.

More explicitly, enumerate the alphabet $XY$ as $z_1,\ldots,z_{an}$, with
$z_1,\ldots,z_n=x_1y_1,\ldots,x_1y_n$. Bump--Gamburd's Proposition~4
(source label \texttt{cpaverage}) gives
\begin{equation}\label{R01:CUE}
 s_r(XY)=\left(\prod_{j=1}^n z_j^d\right)
 \int_{U(d)}\prod_{j=1}^n\det(1+z_j^{-1}U^{-1})
 \prod_{j=n+1}^{an}\det(1+z_jU)\,dU.
\end{equation}
This formula is first read with nonzero $z_j$ and then as a polynomial identity.
Its specialization at identities proves the positive weighted sum
\begin{equation}\label{R01:dimension}
 D_r(an)=\sum_{\mu,\nu\vdash D}
 g(r,\mu,\nu)D_\mu(a)D_\nu(n)
\end{equation}
and the explicit upper bound
\begin{equation}\label{R01:bound}
 0\leq g(\lambda,r,r)\leq
 \left\lfloor\frac{D_r(an)}{D_\lambda(a)}\right\rfloor,
 \qquad
 D_r(an)=\prod_{i=1}^{n}\prod_{j=1}^{d}
 \frac{an+j-i}{n+d-i-j+1}.
\end{equation}
The numerator is a finite rectangular-character moment, not an unknown GCT
coefficient.

\begin{proof}
Apply Schur--Weyl duality to $(\mathbb C^a\otimes\mathbb C^n)^{\otimes D}$
in two ways. Decomposing first for $GL_{an}\times S_D$ gives the factors
$S_\eta(\mathbb C^a\otimes\mathbb C^n)\otimes[\eta]$.
Decomposing separately for $GL_a$ and $GL_n$ gives
$S_\mu(\mathbb C^a)\otimes S_\nu(\mathbb C^n)\otimes
([\mu]\otimes[\nu])$, with diagonal $S_D$ action. Comparing the multiplicity
space of $[r]$ proves \eqref{R01:branching}. Schur orthogonality and
$s_r(B)=\det(B)^d$ prove \eqref{R01:haar}. Equation \eqref{R01:CUE} is the
specialization $N=d,L=n,K=(a-1)n$ of the cited product identity.
Set every $x_i,y_j$ equal to one in \eqref{R01:branching}. Every summand is
nonnegative. Retain the term $(\mu,\nu)=(\lambda,r)$ and use
$D_r(n)=1$; this proves \eqref{R01:bound}. The displayed product is the
hook-content formula in Bump--Gamburd, equation \texttt{e\_hhl} in \S2.2.
\end{proof}

\paragraph{Finite moments and their limiting constants.}
Bump--Gamburd, Corollary~1, gives
\[
 \mathbb E_{U(N)}|\det(1-U)|^{2k}=D_{(N^k)}(2k).
\]
The limit after division by $N^{k^2}$ is
$c_k=\prod_{j=0}^{k-1}j!/(j+k)!$. Thus the finite moment is an integer
dimension, whereas the leading constant generally is not: $c_2=1/12$.
The distinction prevents a scalar asymptotic constant from being treated as
a representation multiplicity.

\subsection{R-02: swaps, tableaux, and positive projectors}

\paragraph{Statement R-02.}
Set $m=an$, $\delta_b=(b-1,b-2,\ldots,0)$, and
$\Delta_b(X)=\det(x_i^{b-j})$. For a tableau $T$ of shape $r$ with
semistandard entries in the ordered alphabet $\{1,\ldots,a\}\times
\{1,\ldots,n\}$, let $p(T)$ and $q(T)$ be the two marginal content vectors.
Write $N_r(p,q)$ for the number of these tableaux with marginal contents
$(p,q)$, and set it equal to zero when a content component is negative.
Then
\begin{align}\label{R02:tableaux}
 g(\lambda,r,r)
 &= [X^{\lambda+\delta_a}Y^{r+\delta_n}]
       \Delta_a(X)\Delta_n(Y)s_r(XY)\\
 &=\sum_{\sigma\in S_a,\tau\in S_n}
    \operatorname{sgn}(\sigma)\operatorname{sgn}(\tau)
    N_r(\lambda+\delta_a-\sigma\delta_a,
        r+\delta_n-\tau\delta_n).\nonumber
\end{align}
The exact unitary sum over swaps acts on the multiplicity by inserting
\begin{equation}\label{R02:swaps}
 s_{(d^n)}(z_1,\ldots,z_m)=
 \sum_{\substack{S\subseteq\{1,\ldots,m\}\\|S|=n}}
 \frac{\prod_{i\in S}z_i^d}
      {\prod_{i\in S,j\notin S}(1-z_j/z_i)}
\end{equation}
in the first line of \eqref{R02:tableaux}, and then setting $z_{ij}=x_i y_j$.
The sum in \eqref{R02:swaps} is combined before specialization or coefficient
extraction; its individual summands have poles and are not separately
nonnegative tableau counts.

There is, nevertheless, an exact positive formula. On the unitary space
\[
 T=S_\lambda(\mathbb C^n\otimes\mathbb C^n)
       \otimes\det(\mathbb C^n)^{-d}\otimes\det(\mathbb C^n)^{-d},
\]
let $P=\int_{U(n)\times U(n)}\rho_T(A,B)\,dA\,dB$. For any orthonormal
basis $(e_i)$ of $T$,
\begin{equation}\label{R02:projector}
 g(\lambda,r,r)=\operatorname{tr}P
   =\|P\|_{\mathrm{HS}}^2=\sum_{i,j}|\langle e_i,Pe_j\rangle|^2.
\end{equation}
In particular, this gives a rigorous sum of squares and a nonvanishing
criterion. It does not give a polynomial-size counting rule: the entries
of $P$ are themselves signed integrals, and the representation dimension
can be exponential in the natural parameters.

\begin{proof}
Multiply \eqref{R01:branching} by the two Vandermonde alternants.
For a partition $\mu$, the alternant $\Delta_a(X)s_\mu(X)$ has exactly one
monomial with strictly decreasing exponent vector equal to
$\lambda+\delta_a$ precisely when $\mu=\lambda$, and its coefficient then
is one. The analogous statement holds for $Y$. This proves the first line
of \eqref{R02:tableaux}. Expanding $s_r$ into semistandard tableaux and
the two alternants into their permutation terms proves the second line.
Formula \eqref{R02:swaps} is the rectangular Weyl-character Laplace expansion
used in Bump--Gamburd, Proposition~4 and \S4. For the last assertion,
Haar invariance gives $P^2=P$ and $P^*=P$; its range is exactly $T^{U(n)\times
U(n)}$. The tensor-subgroup branching identity identifies the dimension of
that range with $g(\lambda,r,r)$. Taking traces of $P^*P=P$ proves
\eqref{R02:projector}.
\end{proof}

\paragraph{What the arithmetic CFKRS recipe supplies.}
CFKRS's unitary formula is a theorem (Theorem~1.5.2 in the supplied source),
and its swap/contour conversion is Lemma~2.5.1. The corresponding
$L$-function mean-value assertion is Conjecture~1.5.1, derived by the
recipe in \S2.1 and the sum over $\Xi$ in \S2.2, equation
\texttt{conj:M}. Replacing the unitary factors in \eqref{R02:swaps} by the
arithmetic Euler product and shifted zeta factors does not define a
character restriction map. Neither RH nor GRH supplies the discarded
off-diagonal estimates in that recipe. The proven action on GCT is
\eqref{R02:tableaux}--\eqref{R02:swaps}; the arithmetic recipe remains a
separate conjectural mean-value statement.

\subsection{R-03: three explicit limits on positivity transfer}

\paragraph{Statement R-03.}
The following three proposed simplifications are false.
First, positivity of a rectangular CUE moment does not imply positivity
of every rectangular Kronecker coefficient. Second, the connected
left-right stabilizer multiplicity need not equal the full determinant
stabilizer multiplicity. Third, one cannot replace tensor-subgroup
branching by one multiplicity-free interlacing step.

\begin{proof}
For $r=(2,2)$ the $S_4$ representation $[r]$ is the two-dimensional
nontrivial summand in the permutation representation on the three
perfect matchings of four letters. Its character at the cycle types
$1^4,21^2,2^2,31,4$ is respectively $2,0,2,-1,0$; the class sizes are
$1,6,3,8,6$. The standard character $[31]$ is $3,1,-1,0,-1$.
Their character inner product gives
\[
 g((3,1),(2,2),(2,2))=(12-12)/24=0,
\]
although $D_{(2,2)}(4)=20>0$. The same elementary calculation yields
\[
 [22]\otimes[22]=[4]\oplus[22]\oplus[1111],\qquad
 \operatorname{Sym}^2[22]=[4]\oplus[22].
\]
Thus $g((1^4),(22),(22))=1$ but the corresponding symmetric Kronecker
coefficient is zero. These values are checked independently by enumerating
the 24 permutations in \texttt{check\_small\_cases.py}.

The tensor homomorphism $U(a)\times U(n)\longrightarrow U(an)$ has kernel
$\{(zI_a,z^{-1}I_n):|z|=1\}$. Its image has one-dimensional connected
center, whereas the block-diagonal Levi $U(a)\times U(n)$ has
two-dimensional connected center. For $a,n\geq2$ these are different
branching problems. Moreover, R-07 below constructs two independent
highest-weight invariants for $g((12,4),(4^4),(4^4))$, so its value is at
least two (in fact exactly two). A single $0/1$ interlacing rule cannot
be its multiplicity formula.
\end{proof}

The $0/1$ Gan--Gross--Prasad periods cited in Alp\"oge's steps~4 and~7
therefore require an additional, explicit comparison of branching
functors before they can control these multiplicities. Iterating
multiplicity-free restrictions can create multiple paths; no
contradiction with such restrictions is asserted. Likewise the squared
Gram formulas are not statements that every Weingarten summand is
positive: Collins--\'Sniady, Corollary~2.4 and its following example, give
$\operatorname{Wg}_N((12))=-1/(N(N^2-1))$ for $N\geq2$.

\subsection{R-04: plethysm requires the joint law of powers}

\paragraph{Statement R-04.}
For $N\geq\ell(\lambda)$ and $|\lambda|=nd$, the plethysm multiplicity is
\begin{equation}\label{R04:plethysm}
 a_\lambda(d[n])=\int_{U(N)}\overline{s_\lambda(U)}
 \sum_{\kappa\vdash d}\frac{1}{z_\kappa}
       \prod_{j\geq1}s_{(n)}(U^j)^{m_j(\kappa)}\,dU,
 \qquad z_\kappa=\prod_j j^{m_j(\kappa)}m_j(\kappa)!.
\end{equation}
Rains's Theorem~1.3 determines the eigenvalue law of each single $U^j$,
but replacing the different powers in \eqref{R04:plethysm} by independent
matrices with those marginal laws is invalid. A quantitative example is
\begin{equation}\label{R04:counter}
 \int_{U(2)}(\operatorname{tr}U)^4
        \overline{(\operatorname{tr}U^2)^2}\,dU=2,
\end{equation}
whereas the product of the two marginal expectations is zero.

\begin{proof}
For any finite-dimensional $W$, the cycle-index formula for the
symmetrizer $d!^{-1}\sum_{\sigma\in S_d}\sigma$ gives
$\chi_{\operatorname{Sym}^dW}(g)=\sum_{\kappa\vdash d}
z_\kappa^{-1}\prod_j\chi_W(g^j)^{m_j(\kappa)}$.
Set $W=\operatorname{Sym}^n\mathbb C^N$ and apply Schur orthogonality.
For \eqref{R04:counter}, in two variables
\[
 p_1^4=s_{(4)}+3s_{(3,1)}+2s_{(2,2)},\qquad
 p_2^2=s_{(4)}-s_{(3,1)}+2s_{(2,2)}.
\]
Orthogonality gives $1-3+4=2$. Multiplication of $U$ by a scalar phase
shows that both separate expectations vanish. The two polynomial
expansions are also verified by the exact computation certificate.
\end{proof}

In representation-ring language $U\mapsto U^j$ yields the Adams
operation on characters, not a homomorphism of nonabelian groups.
Already $\psi^2(V)=\operatorname{Sym}^2V-\bigwedge^2V$ has a negative
Schur coefficient. A positive alternative retaining the complete
coupling is
\begin{equation}\label{R04:wreath}
 a_\lambda(d[n])=\dim[\lambda]^{S_n\wr S_d}
 =\sum_{i,j}\left|\frac{1}{(n!)^d d!}
       \sum_{h\in S_n\wr S_d}\rho_\lambda(h)_{ij}\right|^2.
\end{equation}
Indeed, apply Schur--Weyl duality to $V^{\otimes nd}$ and the orthogonal
projector onto the wreath-product invariants, exactly as in R-02.
The formula is positive and exact; its factorial enumeration cost is
explicit, and it is not a solution of the efficient positive-formula
problem.

\subsection{R-05: Haar covariance computes the orbit-closure ranks}

\paragraph{Statement R-05.}
Let $V=\mathbb C^N$, $W=\operatorname{Sym}^nV$, $0\ne f\in W$, and
$X_f=\overline{GL(V)f}$. Put $E_d=\operatorname{Sym}^dW$ and
$v=f^{\otimes d}\in E_d$. Define
\begin{equation}\label{R05:twirl}
 M_f=\int_{U(N)}|\rho(k)v\rangle\langle\rho(k)v|\,dk.
\end{equation}
Under the Hermitian identification of $E_d$ with its conjugate dual,
$\ker M_f$ is the degree-$d$ ideal of $X_f$, and the image is dual to
$\mathbb C[X_f]_d$. Write the isotypic decomposition
\[
 E_d=\bigoplus_{\lambda\vdash nd} L_\lambda\otimes\mathbb C^{a_\lambda},
 \qquad L_\lambda=S_\lambda V,
 \qquad v_\lambda=\sum_{i=1}^{D_\lambda(N)}e_i\otimes m_i.
\]
Then the $\lambda$ block of $M_f$ is
\begin{equation}\label{R05:block}
 \frac{I_{L_\lambda}}{D_\lambda(N)}\otimes R_\lambda(f),
 \qquad R_\lambda(f)=\sum_i|m_i\rangle\langle m_i|.
\end{equation}
The multiplicity of $L_\lambda^*$ in $\mathbb C[X_f]_d$ is therefore
\begin{equation}\label{R05:rank}
 m_\lambda(f)=\operatorname{rank}R_\lambda(f).
\end{equation}
If $F_\lambda$ is the matrix with columns $m_i$, then, for every $s$,
\begin{equation}\label{R05:minors}
 e_s(R_\lambda(f))=
 \sum_{\substack{I\subseteq\{1,\ldots,a_\lambda\}\\|I|=s}}
 \ \sum_{\substack{J\subseteq\{1,\ldots,D_\lambda(N)\}\\|J|=s}}
       |\det(F_{\lambda,I,J})|^2.
\end{equation}
In particular $m_\lambda(f)\geq s$ if and only if this nonnegative
polynomial is positive. A multiplicity obstruction between two orbit
closures is thus an exact comparison of these Gram ranks, including
the boundary equations, rather than of their scalar traces.

\begin{proof}
For $u\in E_d$,
$\langle u,M_fu\rangle=\int|\langle u,\rho(k)v\rangle|^2\,dk$.
This is zero exactly when the continuous nonnegative integrand is
identically zero. The compact group $U(N)$ is Zariski dense in
$GL_N(\mathbb C)$: a regular function vanishing on it has all derivatives
along its real Lie algebra zero, and the complexification of that Lie
algebra is $\mathfrak{gl}_N(\mathbb C)$; analytic continuation on the
connected complex group gives zero. Consequently the same vanishing
holds on $GL(V)f$ and its closure. A degree-$d$ polynomial on $W$ is
exactly a linear functional on $E_d$ evaluated at $f^{\otimes d}$.
This identifies the kernel with the homogeneous ideal and the image
with the dual quotient.

Off-diagonal isotypic blocks of the twirl vanish by Schur's lemma.
Within $L_\lambda\otimes\mathbb C^{a_\lambda}$, averaging
$|e_i\rangle\langle e_j|$ gives
$\delta_{ij}I_{L_\lambda}/D_\lambda(N)$, proving
\eqref{R05:block}. Thus the orbit span is
$L_\lambda\otimes\operatorname{span}\{m_i\}$, which proves
\eqref{R05:rank}. Finally $R_\lambda=F_\lambda F_\lambda^*$.
The coefficient $e_s$ is the sum of its principal $s$-minors;
Cauchy--Binet gives \eqref{R05:minors}. Since a positive semidefinite
matrix has nonnegative eigenvalues, $e_s>0$ exactly when at least $s$
of those eigenvalues are positive.
\end{proof}

For $f=\det_n$ or the padded permanent, \eqref{R05:twirl} is a Haar
integral over $U(n^2)$. Its entries have degree $nd$ in the matrix
entries and degree $nd$ in their conjugates. Thus the exact
Collins--\'Sniady formula evaluates them. The full, dimension-truncated
Weingarten inverse must be used when tensor degree exceeds group
dimension. The remaining task is to control the ranks of the resulting
structured matrices; knowing that the matrices are positive
semidefinite does not determine which minors vanish.

\subsection{R-06: a Gaussian and integral form of the rank certificate}

\paragraph{Statement R-06.}
Use the tensor realization $E_d\subseteq V^{\otimes D}$ and set
$v=f^{\otimes d}$. For an $N\times N$ matrix of indeterminates $Z$, write
\begin{equation}\label{R06:coefficients}
 Z^{\otimes D}v=\sum_{|\alpha|=D}Z^\alpha c_\alpha,
 \qquad C_f=(c_\alpha)_\alpha.
\end{equation}
There are $N^D$ tensor coordinates and
$\binom{N^2+D-1}{D}$ columns. If $G$ has independent standard complex
Gaussian entries, then
\begin{equation}\label{R06:gaussian}
 \mathbb E|G^{\otimes D}v\rangle\langle G^{\otimes D}v|
   =\sum_{|\alpha|=D}\alpha!\,c_\alpha c_\alpha^*.
\end{equation}
Its image is the orbit span from R-05. It has the same kernel and
the same isotypic ranks as the Haar covariance.

Let $f^\lambda=\dim[\lambda]$ and let
\[
 e_\lambda=\frac{f^\lambda}{D!}
       \sum_{\sigma\in S_D}\chi^\lambda(\sigma)\sigma
\]
act by permutation of tensor factors. Then
\begin{equation}\label{R06:integer}
 A_\lambda(f)=\sum_{\sigma\in S_D}\chi^\lambda(\sigma)\sigma C_f,
 \qquad
 \operatorname{rank}A_\lambda(f)=D_\lambda(N)m_\lambda(f).
\end{equation}
When $f$ is an integral symmetric tensor, $A_\lambda(f)$ has integer
entries. If its tensor entries have absolute value at most $H$, then
every entry of $A_\lambda(f)$ has absolute value at most
$D!f^\lambda N^D H^d$.
The same rank is obtained by replacing the coefficient columns by all
evaluations at the finite grid $Z\in\{0,1,\ldots,D\}^{N^2}$ and then
applying $e_\lambda$. This grid and the coefficient matrix are
explicit but not polynomial-size constructions.
An ordinary polynomial with integer monomial coefficients may require
scaling by $n!$ under the symmetric-tensor embedding to meet the
integrality hypothesis; this nonzero scalar does not change any orbit
span or multiplicity in characteristic zero.

\begin{proof}
For a standard complex Gaussian $z$, polar integration gives
$\mathbb E[z^p\bar z^q]=0$ if $p\ne q$ and $p!$ if $p=q$.
Independence proves \eqref{R06:gaussian}. Its image is
$\operatorname{span}\{c_\alpha\}$. A linear functional kills all
$c_\alpha$ exactly when its evaluation on $Z^{\otimes D}v$ is the zero
polynomial. Since invertible matrices are Zariski dense, that is
equivalent to killing the $GL_N$ orbit span. This proves the image and
kernel assertions without analytic approximation.

The central idempotent $e_\lambda$ projects precisely onto the
$L_\lambda$ isotypic summand by Schur--Weyl duality. Consequently
$e_\lambda C_f$ has rank $D_\lambda(N)m_\lambda(f)$, and multiplication
by the nonzero scalar $D!/f^\lambda$ gives \eqref{R06:integer}.
Symmetric-group characters and permutation matrices have integral
entries. Each tensor coordinate of $Z^{\otimes D}v$ is a sum of $N^D$
terms whose coefficients have magnitude at most $H^d$; hence
$|(c_\alpha)_i|\leq N^D H^d$. The character bound
$|\chi^\lambda(\sigma)|\leq f^\lambda$ and the $D!$ permutations give
the stated estimate. Finally a polynomial of individual degree at
most $D$ vanishing on a Cartesian product of $D+1$ distinct scalars is
zero, by induction on the number of variables. Applying this fact to
every linear functional proves equality of the grid and coefficient
spans.
\end{proof}

This replaces signed Haar expansions by an explicitly positive
coefficient Gram matrix. Cauchy--Binet gives positive sums of squared
integer minors after clearing the indicated denominators. It also
exhibits the exact cost: the tensor model has $N^D$ rows, the
symmetrizer has $D!$ terms, and the displayed grid has
$(D+1)^{N^2}$ points. A claimed efficient positivity algorithm must
compress these data or exploit their ranks; positivity alone does
not achieve that compression.

\subsection{R-07: explicit positive witnesses that survive the boundary}

\paragraph{Statement R-07.}
Let $n\geq2$, $d\geq0$, and $0\leq q\leq d$. For
$\lambda=(nd-q,q)$, define
\[
 p_{2,\ldots,n}(q)=[t^q]\prod_{j=2}^n(1-t^j)^{-1}.
\]
There are explicitly constructed, linearly independent highest-weight
vectors, indexed by partitions of $q$ into parts $2,\ldots,n$, in the
degree-$d$ coordinate ring of the determinant orbit closure. In fact
\begin{equation}\label{R07:positive}
 m_\lambda(\det_n)=a_\lambda(d[n])=p_{2,\ldots,n}(q),
 \qquad
 g(\lambda,(d^n),(d^n))\geq p_{2,\ldots,n}(q).
\end{equation}
The last inequality is equality by Manivel, Corollary~2(1) in \S3.3; it is a specialization of his
stable rectangular Kronecker Theorem~1, restated as
Ikenmeyer--Panova, Theorem~2.1. The construction below gives concrete witnesses for this
known numerical positivity and, simultaneously, proves their survival
on the determinant orbit closure.

More generally, with no restriction $q\leq d$,
\begin{equation}\label{R07:alltwo}
 m_\lambda(\det_n)=a_\lambda(d[n])
 \quad\text{for every }\lambda\vdash nd\text{ with }\ell(\lambda)\leq2.
\end{equation}
Therefore no multiplicity obstruction of length at most two can
separate a padded-permanent orbit closure from the determinant orbit
closure in the same ambient space and degree.

\begin{proof}
Write a binary degree-$n$ form as
$F(s,t)=\sum_{j=0}^n\binom nj b_j s^{n-j}t^j$.
In the coordinate representation $\operatorname{Sym}^n\mathbb C^2$,
the upper-triangular raising operator is
$E b_j=j b_{j-1}$, with $E b_0=0$, and the weight of $b_j$ is
$(n-j,j)$. For $2\leq j\leq n$ put
\begin{equation}\label{R07:central}
 Q_j=b_0^{j-1}\sum_{i=0}^j\binom ji b_i
                  \left(-\frac{b_1}{b_0}\right)^{j-i}.
\end{equation}
This is a polynomial: for $i=0$ the factor $b_i=b_0$ cancels the only
possible negative exponent; all $i\geq1$ terms already have
nonnegative powers of $b_0$. The first two examples are
\[
 Q_2=b_0b_2-b_1^2,\qquad
 Q_3=b_0^2b_3-3b_0b_1b_2+2b_1^3.
\]
Under $\exp(uE)$, $b_i$ goes to
$\sum_{h=0}^i\binom ih u^{i-h}b_h$ and $b_1/b_0$ goes to
$b_1/b_0+u$. The binomial theorem in \eqref{R07:central} shows that
the two translations cancel. Thus $E Q_j=0$. Every term has degree
$j$ in the $b_i$ and weight $((n-1)j,j)$.

For a tuple of nonnegative integers $(a_2,\ldots,a_n)$ satisfying
$\sum_{j=2}^n ja_j=q$, define
\begin{equation}\label{R07:witness}
 P_{\mathbf a}=b_0^{d-q}\prod_{j=2}^n Q_j^{a_j}.
\end{equation}
This is a highest-weight polynomial of degree $d$ and weight
$(nd-q,q)$. Evaluate it on determinant pencils
$F(s,t)=\det(sI_n+tY)$, with $Y$ diagonal. Every monic binary form
is obtained in this way over $\mathbb C$ by factoring the corresponding
one-variable polynomial. In particular, one may set $b_0=1,b_1=0$
and choose $b_2,\ldots,b_n$ arbitrarily. On this family $Q_j=b_j$,
so the restrictions of the $P_{\mathbf a}$ are distinct monomials
in independent variables. They are linearly independent.

For completeness, the ambient highest-weight dimension in this range
is obtained by a short weight count. In
$\operatorname{Sym}^d(\operatorname{Sym}^n\mathbb C^2)$ the weight
$(nd-q,q)$ has dimension equal to the number of partitions of $q$
into at most $d$ parts, each at most $n$. When $q\leq d$ this is
$[t^q]\prod_{j=1}^n(1-t^j)^{-1}$. A $GL_2$ irreducible of highest
weight $(nd-r,r)$ has one vector at each weight $(nd-q,q)$ with
$r\leq q\leq nd-r$. Subtracting consecutive weight dimensions
therefore gives $p_{2,\ldots,n}(q)$. This proves that the displayed
independent witnesses form a basis and establishes the first equality
in \eqref{R07:positive} once boundary survival is checked.

Now let $V$ have dimension $n^2$ and view a binary form as an element
of $\operatorname{Sym}^nV^*$ supported on the first two coordinates.
It factors as $\prod_{i=1}^n\ell_i$ and hence equals
$\det(\operatorname{diag}(\ell_1,\ldots,\ell_n))$.
This is a linear, possibly singular, substitution into $\det_n$.
Since $GL(V)$ is dense in $\operatorname{End}(V)$ and substitution
is polynomial, every such form lies in the determinant orbit closure.
A vector of highest weight $\lambda=(\lambda_1,\lambda_2,0,\ldots,0)$
in $\operatorname{Sym}^d(\operatorname{Sym}^nV)$ uses only the first
two basis vectors: each other weight coordinate is a sum of
nonnegative exponents and equals zero. The additional raising
operators annihilate it, so its highest-weight condition is exactly
the binary one. A nonzero polynomial in these binary coordinates
cannot vanish on every binary form. Thus the homogeneous ideal of
the determinant orbit closure has no nonzero highest-weight vector
of length at most two. Taking the quotient proves
\eqref{R07:alltwo}, including all degrees and all two-row weights.

The determinant-coordinate multiplicity is at most its orbit
multiplicity, which is the symmetric rectangular Kronecker coefficient,
and that is at most $g(\lambda,r,r)$; this is the cited stabilizer
statement in B\"urgisser--Ikenmeyer--Panova, \S1(b).
This proves the inequality in \eqref{R07:positive}. Finally any
padded-permanent coordinate ring is a quotient of the same ambient
$\operatorname{Sym}(\operatorname{Sym}^nV)$, so its multiplicity is
at most $a_\lambda(d[n])=m_\lambda(\det_n)$ for two-row $\lambda$.
\end{proof}

For $n=d=q=4$, the two witnesses are $Q_4$ and $Q_2^2$, with
\[
 Q_4=b_0^3b_4-4b_0^2b_1b_3+6b_0b_1^2b_2-3b_1^4.
\]
They specialize to $b_4$ and $b_2^2$, proving directly that the
determinant coordinate multiplicity at $(12,4)$ is two. This is
stronger information than nonnegativity of a Haar integral.

\subsection{R-08: the exact Molien numerator needed for orbit closures}

\paragraph{Statement R-08.}
Let $S=\operatorname{Sym}(W)$ be a polynomial coordinate algebra,
$A=S/I$ a graded $GL_N$-equivariant quotient, and $K=U(N)$. Then
\begin{equation}\label{R08:molien}
 \sum_{d\geq0}m_\lambda(A_d)t^d
   =\int_K\overline{\chi_\lambda(k)}H_A(t,k)\,dk.
\end{equation}
For the ambient algebra,
$H_S(t,k)=\det(1-t\rho_W(k))^{-1}$.
For a finite graded equivariant free resolution
$\cdots\to S\otimes T_1\to S\otimes T_0\to A\to0$,
\begin{equation}\label{R08:numerator}
 H_A(t,k)=
 \frac{\sum_i(-1)^i\chi_{T_i}(t,k)}
      {\det(1-t\rho_W(k))}.
\end{equation}
Thus the Molien denominator and ratios identities compute the desired
multiplicities only after the orbit ideal's equivariant numerator,
or equivalently the ranks in R-05--R-06, is supplied.

\begin{proof}
Coefficientwise Schur orthogonality proves \eqref{R08:molien}.
The eigenvalue expansion of the symmetric algebra gives the
denominator formula. Alternating graded characters are additive
on exact complexes, yielding \eqref{R08:numerator}. For a purely
formal interpretation no convergence assumption is needed; for
analytic evaluation take $|t|<1$ in the unitary polynomial case.
\end{proof}

\subsection{R-09: an exact ratios kernel for squared GCT multiplicities}

\paragraph{Statement R-09.}
Let $n,d,N\geq1$, $r=(d^n)$, and define
\[
 S_N(n,d)=\sum_{\substack{\lambda\vdash nd\\\ell(\lambda)\leq N}}
                    g(\lambda,r,r)^2.
\]
Set $S_0(n,d)=0$.
For $0<t<1$ and $A,B,C,D\in U(n)$ put
\begin{align*}
 \mathcal R_N(t;A,B,C,D)=\int_{U(N)}
 &\det\bigl(I-t(A\otimes B)\otimes U\bigr)^{-1}\\[-2mm]
 &\quad\cdot
 \det\bigl(I-t(C\otimes D)\otimes U^{-1}\bigr)^{-1}\,dU.
\end{align*}
Then the explicitly specified fourfold Fourier projection is
\begin{equation}\label{R09:squares}
 \int_{U(n)^4}\mathcal R_N(t;A,B,C,D)
   (\det A\det B\det C\det D)^{-d}\,dA\,dB\,dC\,dD
       =t^{2nd}S_N(n,d).
\end{equation}
For $N\geq n^2$, the inner ratio is evaluated exactly by
Bump--Gamburd, Theorem~3, at $K=L=0,Q=R=n^2$:
\begin{equation}\label{R09:evaluated}
 \mathcal R_N(t;A,B,C,D)
       =\det\bigl(I_{n^4}-t^2 A\otimes B\otimes C\otimes D\bigr)^{-1}.
\end{equation}
Consequently an explicitly evaluated Molien--Weyl coefficient is the
sum of squares of precisely the rectangular Kronecker coefficients
needed in GCT. The truncated sums satisfy
\begin{equation}\label{R09:monotone}
 S_N(n,d)-S_{N-1}(n,d)
    =\sum_{\substack{\lambda\vdash nd\\\ell(\lambda)=N}}g(\lambda,r,r)^2
    \geq0,
 \qquad
 g(\lambda,r,r)\leq\lfloor\sqrt{S_N(n,d)}\rfloor
 \quad(\ell(\lambda)\leq N).
\end{equation}
The stabilized value has a second nonnegative moment expression:
\begin{equation}\label{R09:finitefourth}
 S_{n^2}(n,d)=\frac1{(nd)!}\sum_{\sigma\in S_{nd}}\chi^r(\sigma)^4.
\end{equation}

\begin{proof}
The Cauchy identity and absolute convergence for $t<1$ expand the two
inverse determinants into
$\sum_\lambda t^{|\lambda|}s_\lambda(A\otimes B)s_\lambda(U)$
and
$\sum_\mu t^{|\mu|}s_\mu(C\otimes D)s_\mu(U^{-1})$.
Schur orthogonality on $U(N)$ retains exactly $\lambda=\mu$ of
length at most $N$. The two applications of the tensor-subgroup
projection in R-02 then retain $|\lambda|=nd$ and give
$g(\lambda,r,r)$ from each pair of $U(n)$ factors. This proves
\eqref{R09:squares}. When $N\geq n^2$, the length cutoff is
inactive because both product alphabets have $n^2$ entries. A second
Cauchy identity gives \eqref{R09:evaluated}, equivalently the stated
specialization of the ratios theorem. The definition immediately
proves \eqref{R09:monotone}. Finally the character of
$[r]\otimes[r]$ is $(\chi^r)^2$. Its squared character norm is
the sum of squares of its irreducible multiplicities, proving
\eqref{R09:finitefourth}; symmetric-group characters are real.
The tensor-subgroup branching formula also gives
$g(\lambda,r,r)=0$ for $\ell(\lambda)>n^2$, because
$S_\lambda(\mathbb C^n\otimes\mathbb C^n)=0$ in that range.
Thus the cutoff in the left side of \eqref{R09:finitefourth}
omits no nonzero multiplicity.
\end{proof}

The positive quantities in \eqref{R09:squares} are obtained after
projection. The pointwise ratio and individual swap summands need
not be nonnegative. An approximation argument would need a
certified error below one half to recover an integer coefficient,
with bounds uniform in the growing tensor degree. Neither the
closed ratio formula nor RH supplies those estimates or compresses
the remaining fourfold Fourier projection.

\subsection{R-10: the complete height-two rectangle and its boundary gate}

\paragraph{Statement R-10 (known formula, specialized and compared).}
Let $d\geq0$ and $\lambda\vdash2d$. For partitions of length at
most four, append zero parts to obtain $(\lambda_1,\ldots,\lambda_4)$.
Then
\begin{equation}\label{R10:rectangle}
g(\lambda,(d,d),(d,d))=
\begin{cases}
1,&\ell(\lambda)\leq4\text{ and all four parts have the same parity},\\
0,&\text{otherwise}.
\end{cases}
\end{equation}
More precisely, the multiplicity of $[\lambda]$ in
$\operatorname{Sym}^2[d,d]$ is one exactly for partitions with all
parts even and length at most four. Its multiplicity in
$\bigwedge^2[d,d]$ is one exactly for partitions with exactly four
odd parts. These formulas are the constant-term specialization of
Blasiak--Mulmuley--Sohoni, GCT~IV, Corollary~18.4
(source label \texttt{c kronecker generating function sym ext});
they are cited known results, not new numerical formulas.

In the determinant problem with $n=2$, $N=4$, the boundary side is
also completely determined:
\begin{equation}\label{R10:closure}
m_\lambda(\det_2)=a_\lambda(d[2])=
\begin{cases}
1,&\ell(\lambda)\leq4\text{ and every part of }\lambda\text{ is even},\\
0,&\text{otherwise}.
\end{cases}
\end{equation}
Thus all four-odd-part summands of the ordinary rectangular
coefficient are absent from the determinant coordinate ring.
For $d=2$, this includes the $\lambda=(1,1,1,1)$ example in R-03.
Simultaneous conjugation of the two inner partitions also extends
\eqref{R10:rectangle}, as an ordinary Kronecker formula, to rectangles
of width two.

\begin{proof}
Write $c_i=\lambda_i-\lambda_{i+1}$ for $1\leq i\leq3$ and
$c_4=\lambda_4$, so $c_i$ is the number of columns of height $i$.
GCT~IV, Corollary~18.4, gives the generating function
$\sum_{\alpha\vdash2d,\ell(\alpha)\leq2}
g_{\varepsilon,\alpha,\lambda}x^{\alpha_1-\alpha_2}$.
Three of its four cases have a factor $x$ or $x^2$ multiplying
polynomials with nonnegative exponents, and hence contribute no
constant term. The remaining case is the product
$\llbracket c_1\rrbracket\llbracket c_2\rrbracket
\llbracket c_3\rrbracket$, where
$\llbracket b\rrbracket=x^b+x^{b-2}+\cdots$ ends at $1$ or $x$.
Its constant coefficient is one exactly when $c_1,c_2,c_3$ are all
even; the case condition then requires
$\varepsilon=(-1)^{c_3+c_4}=(-1)^{\lambda_4}$.
Thus the symmetric square retains the all-even partitions and the
exterior square retains the four-odd-part partitions. Their sum is
\eqref{R10:rectangle}. Length greater than four gives zero by
tensor-subgroup branching. Finally
$[\rho']=[\rho]\otimes\operatorname{sgn}$, so conjugating both
inner partitions leaves their tensor product unchanged.

The determinant $x_{11}x_{22}-x_{12}x_{21}$ is a nondegenerate
quadratic form in four variables. Over $\mathbb C$ all nondegenerate
quadratic forms of a fixed dimension are congruent. They form a
nonempty Zariski-open subset of the space of quadrics, so
$\operatorname{Det}_2=\operatorname{Sym}^2(\mathbb C^4)^*$.
Consequently its coordinate ring is
$\operatorname{Sym}(\operatorname{Sym}^2\mathbb C^4)$.

Here is a direct highest-weight proof of its familiar multiplicity
formula. In $\operatorname{Sym}(\operatorname{Sym}^2\mathbb C^s)$
write the generators as the entries of a symmetric matrix $Q$,
with $q_{ij}$ of weight $e_i+e_j$. The upper-triangular action on
these generators is $Q\mapsto g^tQg$.
Let $\Delta_i$ be the leading principal $i\times i$ minor and
$\Delta_0=1$. Each $\Delta_i$ is fixed by the upper unitriangular
group and has weight $2(e_1+\cdots+e_i)$.
On the dense open set where all $\Delta_i$ are nonzero, elementary
upper-triangular congruences reduce $Q$ to the unique diagonal
matrix with entries $\Delta_i/\Delta_{i-1}$.
Therefore a nonzero unitriangular-invariant polynomial has nonzero
restriction to diagonal matrices. If it is also a torus eigenvector
of weight $\lambda$, this restriction is a scalar multiple of
$\prod_i q_{ii}^{\lambda_i/2}$. It exists only when each
$\lambda_i$ is even, and its value determines the polynomial on a
dense open set. Hence each highest-weight space has dimension at
most one and vanishes unless $\lambda=2\mu$.
Conversely, for every partition $\mu$ of length at most $s$,
\[
 \prod_{i=1}^s\Delta_i^{\mu_i-\mu_{i+1}},\qquad\mu_{s+1}=0,
\]
is a nonzero highest-weight polynomial of weight $2\mu$ and
degree $|\mu|$. This proves
$\operatorname{Sym}^d(\operatorname{Sym}^2\mathbb C^s)
=\bigoplus_{\mu\vdash d,\ell(\mu)\leq s}S_{2\mu}\mathbb C^s$
and, at $s=4$, \eqref{R10:closure}.
\end{proof}

The stabilizer interpretation agrees with these formulas.
The full linear stabilizer of the nondegenerate quadratic form is
$O_4$, whereas the connected tensor image of $SL_2\times SL_2$ is
$SO_4$: its kernel is $\{(I,I),(-I,-I)\}$, so its dimension is six,
the dimension of $SO_4$. The extra orthogonal component removes
the odd sector. Thus the connected-stabilizer upper bound, the
symmetric rectangular coefficient, and the closure multiplicity
are completely distinguished in this case.

\subsection{R-11: small-degree padding excludes all four-row obstructions}

\paragraph{Statement R-11.}
Work over $\mathbb C$. Let $1\leq m\leq3$, $n\geq m$, $N=n^2$,
$V=\mathbb C^N$, $f\in\operatorname{Sym}^mV^*$, and $\ell\in V^*$.
Set $h=\ell^{n-m}f$. For every coordinate-ring degree $d\geq0$
and every $\lambda\vdash nd$ with $\ell(\lambda)\leq4$,
\begin{equation}\label{R11:fourrow}
 m_\lambda(h)\leq m_\lambda(\det_n).
\end{equation}
The assertion includes padded permanents of degree at most three.
It is a deduction from classical linear determinantal
representations, not a claim that arbitrary $f$ itself has an
$n\times n$ determinantal representation.

\begin{proof}
First consider forms in at most four variables. Linear forms are
$1\times1$ determinants. Every quadric in at most four variables
is congruent over $\mathbb C$ to $y_1^2+\cdots+y_r^2$, $r\leq4$.
With unused $y_i$ set to zero it is the determinant
\begin{equation}\label{R11:quadric}
 \det\begin{pmatrix}
 y_1+\mathrm i y_2&y_3+\mathrm i y_4\\
 -y_3+\mathrm i y_4&y_1-\mathrm i y_2
 \end{pmatrix}=
 y_1^2+y_2^2+y_3^2+y_4^2.
\end{equation}
Beauville (2000), Corollary~6.4, states over an algebraically
closed field that every smooth cubic surface has a $3\times3$
linear determinantal equation. Thus over $\mathbb C$ every smooth
quaternary cubic is a $3\times3$ linear determinant, after
absorbing its nonzero scalar factor in one matrix row.
Smooth cubics are Zariski dense in the affine space of quaternary
cubics: the discriminant-complement is open and nonempty, as the
Fermat cubic shows. Therefore every quaternary cubic lies in the
closure of the linear-determinant image.

If $q=\det M$ has degree $m$ and $u$ is linear, then
$u^{n-m}q=\det\operatorname{diag}(M,uI_{n-m})$.
The entries of this block matrix give a linear substitution into
$\det_n$. Such substitutions belong to $\operatorname{Det}_n$,
because invertible linear maps are dense in all linear maps and
substitution is polynomial. For fixed $u$, the preimage of
$\operatorname{Det}_n$ under the linear map $q\mapsto u^{n-m}q$
is closed. The density just proved extends the assertion from
smooth cubics to every cubic. We have shown that
\begin{equation}\label{R11:projectedcontainment}
 u^{n-m}q\in\operatorname{Det}_n.
\end{equation}
Here $q$ has degree $m\leq3$, and $q$ and $u$ are supported
on the same space of at most four variables.

Put $k=\min(4,N)$ and let $\pi$ set the last $N-k$ variables
equal to zero. Let $P$ be a highest-weight vector of weight
$\lambda$ in the degree-$d$ ideal of $\operatorname{Det}_n$.
Since $\lambda_j=0$ for $j>k$ and all coordinate generators have
nonnegative weights, every monomial of $P$ uses only coefficients
of forms supported on the first $k$ variables. Hence
$P(F)=P(\pi F)$ for every $F\in\operatorname{Sym}^nV^*$.
For $g\in GL(V)$ the projection is multiplicative on forms, so
\[
 \pi(g\cdot h)=(\pi(g\cdot\ell))^{n-m}\pi(g\cdot f)
 \in\operatorname{Det}_n
\]
by \eqref{R11:projectedcontainment}. It follows that
$P(g\cdot h)=0$ for every $g$, and hence $P$ belongs to the
ideal of $\overline{GL(V)\cdot h}$.
Thus the highest-weight space of the determinant ideal is
contained in that of the padded-form ideal for every such
$\lambda$. Finite-dimensional polynomial $GL(V)$-modules are
semisimple in characteristic zero. Subtracting the ideal
multiplicity from the common ambient multiplicity gives
\eqref{R11:fourrow}. If $N<4$, partitions of length greater than
$N$ already have multiplicity zero, so the same statement holds.
\end{proof}

\subsection{R-12: exact permanent-pencil dominance and four-row deficits}

\paragraph{Statement R-12 (computer-assisted finite-degree deduction).}
For $n=4,5$, the polynomial map
\[
 \Phi_n:(\operatorname{Mat}_n\mathbb C)^4
 \longrightarrow\operatorname{Sym}^n(\mathbb C^4)^*,
 \qquad(A_1,A_2,A_3,A_4)\longmapsto
 \operatorname{per}\!\left(\sum_{k=1}^4x_kA_k\right)
\]
is dominant. Consequently, for every $d\geq0$ and
$\lambda\vdash nd$ with $\ell(\lambda)\leq4$,
\begin{equation}\label{R12:perambient}
 m_\lambda(\operatorname{per}_n)=a_\lambda(d[n])
 \qquad(n=4,5).
\end{equation}
Combining the quartic case with G-07 and writing $D=320112$ gives
the exact comparison
\begin{equation}\label{R12:deficit}
 m_\lambda(\operatorname{per}_4)-m_\lambda(\det_4)
 =a_{\lambda-(D,D,D,D)}((d-D)[4]),
\end{equation}
with the same zero conventions as in G-07.
In particular, for every integer $s\geq0$,
\begin{equation}\label{R12:gapone}
 d=D+s,\qquad \lambda=(D+4s,D,D,D)
 \quad\Longrightarrow\quad
 m_\lambda(\operatorname{per}_4)-m_\lambda(\det_4)=1.
\end{equation}
These are multiplicity obstructions; no occurrence-vanishing
claim is implied. For $n=5$, G-06 and \eqref{R12:perambient}
give a four-row multiplicity obstruction in some degree, without
specifying that degree or highest weight.

\paragraph{Explicit rank certificate for $n=4$.}
Each entry of the following array is the coefficient quadruple of
a linear form in the ordered variables $(x_1,x_2,x_3,x_4)$:
\begin{equation}\label{R12:pencil4}
\begin{pmatrix}
(2,-1,0,0)&(2,-2,0,-2)&(-2,-1,-1,0)&(0,-1,-1,2)\\
(0,1,0,0)&(2,-2,-1,1)&(-2,1,-2,0)&(-1,-2,-1,-1)\\
(0,-2,0,2)&(1,-2,-1,-1)&(2,-1,-1,0)&(2,2,0,0)\\
(1,1,2,1)&(1,0,2,-2)&(-1,2,-1,1)&(2,1,1,1)
\end{pmatrix}.
\end{equation}
Call the resulting matrix of linear forms $L_4(x)$. Index the
parameters $a_{ijk}$ by $0\leq i,j\leq3$, $0\leq k\leq3$,
where $a_{ijk}$ multiplies $x_{k+1}$ in entry $(i+1,j+1)$,
and assign column number $c=4(4i+j)+k$.
Order the 35 exponent quadruples $\alpha$ of total degree four in
decreasing lexicographic order. Define the integer Jacobian by
\begin{equation}\label{R12:derivative}
 J_{\alpha,4(4i+j)+k}
 =[x^\alpha]\;x_{k+1}
 \operatorname{per}\bigl(L_4(x)_{\widehat{i+1},\widehat{j+1}}\bigr).
\end{equation}
Take all 35 rows and the ordered columns
\[
 C_4=(0,1,\ldots,30,32,33,34,35).
\]
The resulting square integer matrix has determinant
\begin{equation}\label{R12:minor4}
 \det J_{C_4}
 =45459683955923995466096193465479594741071872
 \equiv77\pmod{101}.
\end{equation}
Thus \eqref{R12:pencil4}--\eqref{R12:minor4} specify the entire
nonvanishing certificate in the text, independently of a file format.

\paragraph{Explicit rank certificate for $n=5$.}
Use the same coefficient-quadruple convention with the array
\begin{equation}\label{R12:pencil5}
\begin{pmatrix}
(2,2,-1,1)&(2,-2,-2,1)&(-2,-1,1,1)&(1,-1,0,0)&(1,2,0,2)\\
(1,0,0,-1)&(1,2,0,2)&(1,2,-1,2)&(0,2,-1,1)&(-1,2,1,0)\\
(-2,-2,1,0)&(1,-1,-1,2)&(0,1,0,-1)&(2,2,-1,1)&(2,-1,1,-2)\\
(2,0,0,1)&(-1,1,2,2)&(0,0,1,-2)&(1,-1,1,-1)&(-1,0,1,-2)\\
(0,-2,1,-1)&(1,0,-1,1)&(-1,-1,1,-1)&(-1,-1,2,1)&(0,-2,-2,2)
\end{pmatrix}.
\end{equation}
There are 56 degree-five coefficient rows, again in decreasing
lexicographic order, and parameter column $4(5i+j)+k$.
Use the same derivative rule, with $4\times4$ complementary
permanents, and the columns
\[
 C_5=(0,1,\ldots,37,39,40,\ldots,56).
\]
The selected $56\times56$ minor has determinant $30$ modulo 101.

\begin{proof}
Differentiating the permanent with respect to one matrix entry
gives the complementary permanent; differentiating that entry
with respect to $a_{ijk}$ contributes $x_{k+1}$. This proves the
Jacobian rule. Expanding the complementary permanents and
performing exact modular elimination gives the displayed
nonzero residues. For full reproducibility,
\texttt{check\_permanent\_projection.py} records the full pencils,
coefficient matrices, selected columns, and integer minors.
The characteristic-zero determinants are also calculated by
fraction-free Bareiss elimination. As a second check on the
derivative expansion, every one of the 64 quartic columns was
recomputed as the exact finite difference
$\Phi_4(a+e_{ijk})-\Phi_4(a)$; this is exactly the derivative
because the map is affine in each individual parameter.
Twelve quintic columns were checked in the same way.

A nonzero minor modulo 101 is a nonzero integer minor, so the
complex differential is surjective at each displayed point.
Fix the other parameters and retain the selected 35 or 56
parameters. The resulting square polynomial map has invertible
complex differential there. By the inverse function theorem its
image contains a nonempty analytic open subset. A polynomial
vanishing on that image vanishes on an analytic open subset and
is identically zero. This proves Zariski density, hence dominance.

Every permanent pencil is a linear endomorphism substitution
into $\operatorname{per}_n$, and therefore belongs to its
$GL_{n^2}$ orbit closure. This closed variety consequently
contains the full linear space of degree-$n$ forms supported on
the first four variables. A highest-weight polynomial of weight
$\lambda$ with $\ell(\lambda)\leq4$ uses only those coefficient
variables, by the nonnegative-weight support argument of R-07.
If it vanishes on the permanent orbit closure, it vanishes on
that entire coefficient space, so it is the zero polynomial.
Thus the ideal has no highest-weight vector in those weights,
which proves \eqref{R12:perambient}.

Equation \eqref{R12:deficit} now follows by subtracting G-07.
For \eqref{R12:gapone}, the shifted partition is $(4s)$, whose
multiplicity in $\operatorname{Sym}^s(\operatorname{Sym}^4V)$ is
one: its weight space is spanned by $(e_1^4)^s$, and this vector
is highest weight. This also covers $s=0$.
For $n=5$, the four-variable determinant image has dimension at
most $52<56$ by G-06, so its ideal contains a nonzero homogeneous
highest-weight vector. G-05 excludes highest weights of length
at most three. Hence it removes at least one four-row ambient
multiplicity, whereas \eqref{R12:perambient} retains them all.
\end{proof}

These computations are exact finite certificates. The independent
agent recheck uses Ryser's formula at integer evaluation points
and an invertible interpolation matrix, so it does not reuse
the complementary-permanent coefficient-expansion algorithm.
The checks establish the stated degrees four and five; they
do not extrapolate dominance to growing $n$ or supply an
asymptotic permanent-versus-determinant lower bound.

\subsection{R-13: the exact range of quaternary permanent dominance}

\paragraph{Statement R-13 (finite certificates and a dimension obstruction).}
For the polynomial map $\Phi_n$ in R-12,
\begin{equation}\label{R13:range}
 \Phi_n\text{ is dominant over }\mathbb C
 \quad\Longleftrightarrow\quad1\leq n\leq16.
\end{equation}
Consequently, for $2\leq n\leq16$ and every
$\lambda\vdash nd$ with $\ell(\lambda)\leq4$,
\begin{equation}\label{R13:multiplicities}
 m_\lambda(\operatorname{per}_n)=a_\lambda(d[n]).
\end{equation}
For every $4\leq n\leq16$, the unpadded permanent and determinant
of size $n$ have a multiplicity obstruction with exactly four rows
in some coordinate-ring degree. R-12 gives an explicit infinite
family when $n=4$; no degree or partition is asserted here for
$5\leq n\leq16$. This is a certificate-based finite-range result,
with no claim of literature priority or asymptotic separation.

\paragraph{Complete specification of the finite certificates.}
All arithmetic for this paragraph is in $\mathbb F_{101}$, with
integer coefficient representatives in $\{0,\ldots,100\}$.
For each $n=1,\ldots,16$ set
\[
\begin{gathered}
 s_0=20261121+1009n,\qquad
 s_{c+1}=(1664525s_c+1013904223)\bmod2^{32},\\
 a_c=\left\lfloor\frac{s_{c+1}}{2^{16}}\right\rfloor\bmod101.
\end{gathered}
\]
There are $4n^2$ coefficients, with $0\leq c<4n^2$, and the
pencil is
\[
 L_n(x)_{i+1,j+1}=\sum_{k=0}^3a_{4(ni+j)+k}x_{k+1}
 \quad(0\leq i,j<n).
\]
Put $b_n=\binom{n+3}{3}$ and order all integer points
$\alpha\in\mathbb N^4$ with $|\alpha|=n$ in decreasing
lexicographic order. Define the $b_n\times4n^2$ matrix
\begin{equation}\label{R13:evaluationmatrix}
 E_n\bigl(\alpha,4(ni+j)+k\bigr)
 =\alpha_{k+1}\operatorname{per}
   \bigl(L_n(\alpha)_{\widehat{i+1},\widehat{j+1}}\bigr)
 \pmod{101}.
\end{equation}
Choose columns greedily from left to right, retaining a column
exactly when it increases the span of previously retained columns.
Equivalently, scan columns left to right in Gaussian elimination,
choose the first available nonzero pivot row, and clear the rows
below it. This produces $b_n$ columns $C_n$ in every case, and the
square minors, with rows and columns in the stated order, have
the following nonzero residues:
\[
\begin{array}{c|r|r@{\qquad}c|r|r}
n&b_n&\det(E_n)_{C_n}&n&b_n&\det(E_n)_{C_n}\\\hline
1&4&1&9&220&67\\
2&10&11&10&286&67\\
3&20&64&11&364&42\\
4&35&49&12&455&87\\
5&56&85&13&560&25\\
6&84&28&14&680&69\\
7&120&25&15&816&73\\
8&165&59&16&969&8
\end{array}
\]
These uniform range certificates use the displayed generator;
R-12 retains its separate, smaller-entry quartic and quintic pencils.
The directory \texttt{permanent\_range} contains all generated
integer coefficients, evaluation points, chosen columns, and full
matrices, with unambiguous row-major byte encoding and checksums.
Thus the formula above and the code reproduce the certificates
without probabilistic choices or an external data source.

\begin{proof}
For the positive direction, \eqref{R13:evaluationmatrix} is the
derivative of the coefficient map followed by evaluation at the
specified points, by the complementary-permanent derivative
rule in R-12. Its full row rank $b_n$ forces the coefficient-map
differential, whose target dimension is $b_n$, to be surjective.
The nonzero minor modulo 101 is the reduction of a nonzero
integer minor. Hence the characteristic-zero differential has
the same required full rank, and the inverse-function argument
in R-12 proves dominance.

Here are the exact algorithms used to certify every entry and
minor in the table. For a numeric $n\times n$ matrix $M$, let
$F(\varnothing)=1$ and
\[
 F(S)=\sum_{j\in S}M_{|S|,j}F(S\setminus\{j\}).
\]
Induction on $|S|$ shows that $F(S)$ is the permanent formed from
the first $|S|$ rows and columns $S$. Thus $F([n])$ is the
permanent. Reverse differentiation of this recurrence computes
every entry derivative, hence every complementary permanent,
in $O(n2^n)$ field operations per evaluation point. This is
implemented in \texttt{check\_permanent\_range.cpp}.

All $3{,}348{,}864$ entries across the sixteen matrices were then
recomputed from Ryser's different identity
\[
 \operatorname{per}(M)=
 \sum_{S\subseteq[n]}(-1)^{n-|S|}
 \prod_{i=1}^n\left(\sum_{j\in S}M_{ij}\right).
\]
To prove this identity, expand the products as maps
$f:[n]\to[n]$. The coefficient of a fixed map is
$\sum_{S\supseteq\operatorname{im}f}(-1)^{n-|S|}$, which is zero
unless $f$ is onto and is one if it is onto. Such onto maps are
exactly permutations. Differentiating this identity and using
Gray-code row-sum updates gives the independent algorithm in
\texttt{check\_permanent\_range\_ryser.cpp}. Every entry agrees
with the first computation. Finally
\texttt{check\_permanent\_range\_minors.py} recomputes each
selected determinant by a second modular elimination algorithm;
all sixteen nonzero residues agree. These are finite exact
checks, rather than floating-point rank estimates.

For the converse, the torus
\[
 H_n=\{(u_1,\ldots,u_n,v_1,\ldots,v_n)\in(\mathbb C^*)^{2n}:
                  \prod_i u_i\prod_j v_j=1\}
\]
preserves $\Phi_n$ under the entrywise action
$L_{ij}\mapsto u_i v_jL_{ij}$. It has dimension $2n-1$.
On the nonempty open set where every entry $L_{ij}$ is a nonzero
linear form, a stabilizing torus element satisfies
$u_i v_j=1$ for all $i,j$. It is therefore of the form
$u_i=t$, $v_j=t^{-1}$, a one-dimensional stabilizer.
Thus each generic fiber contains an orbit of dimension $2n-2$,
and
\begin{equation}\label{R13:dimension}
 \dim\overline{\operatorname{im}\Phi_n}
       \leq4n^2-2n+2.
\end{equation}
At $n=17$, this is $1124<1140=b_{17}$. For $n\geq18$,
\[
 6\bigl(b_n-(4n^2-2n+2)\bigr)
 =n^3-18n^2+23n-6>0.
\]
Consequently no $n\geq17$ map is dominant. This proves
\eqref{R13:range}.

For $n\geq2$, the ambient variable space has dimension $n^2\geq4$,
so dominance places the entire space of degree-$n$ forms in the
first four variables inside the permanent orbit closure. The
highest-weight support argument of R-12 proves
\eqref{R13:multiplicities}.
For $n\geq4$, G-06 gives a nonzero determinant equation supported
on four variables, and G-05 excludes weights with at most three
rows. Thus some exactly four-row highest weight has determinant
multiplicity below its ambient value. Equation
\eqref{R13:multiplicities} retains that ambient value on the
permanent side when $n\leq16$, giving the asserted obstruction.
\end{proof}

The dimension ceiling concerns full ambient survival in four
variables. It does not determine the comparison between the two
proper four-variable images for $n\geq17$.

\subsection{R-14: a dimension criterion that survives genuine padding}

\paragraph{Statement R-14.}
Let $\mathcal P_{m,4}$ be the Zariski closure of the image of
$\Phi_m$ and put $c_m=\dim\mathcal P_{m,4}$.
For integers $n>m\geq2$, embed the usual padded permanent
$h_{m,n}=z^{n-m}\operatorname{per}_m$ in
$\operatorname{Sym}^n(\mathbb C^{n^2})^*$, with $z$ independent
of the $m^2$ permanent variables. The projection of its orbit
closure to the first four variables is exactly
\begin{equation}\label{R14:paddingvariety}
 \mathcal Z_{m,n}=
 \{\ell^{n-m}q:\ell\in(\mathbb C^4)^*,\ q\in\mathcal P_{m,4}\},
 \qquad \dim\mathcal Z_{m,n}=c_m+3.
\end{equation}
If
\begin{equation}\label{R14:criterion}
 c_m+3>2n^2+2,
\end{equation}
then, for every sufficiently large coordinate-ring degree $d$,
there is a partition $\lambda\vdash nd$ of length exactly four
such that
\[
 m_\lambda(h_{m,n})>m_\lambda(\det_n).
\]
No effective threshold for $d$ is asserted.

For $4\leq m\leq16$, R-13 gives $c_m=\binom{m+3}{3}$, so the
sufficient condition becomes
\begin{equation}\label{R14:explicitrange}
 m<n\leq\left\lfloor\sqrt{\frac12\binom{m+3}{3}}\right\rfloor.
\end{equation}
The nonempty padded ranges are
\[
\begin{array}{c|ccccccccc}
m&8&9&10&11&12&13&14&15&16\\\hline
n&9&10&11&12\text{--}13&13\text{--}15&14\text{--}16&
15\text{--}18&16\text{--}20&17\text{--}22
\end{array}.
\]
In particular, some four-row multiplicity obstruction separates
the unpadded $\operatorname{per}_n$ and $\det_n$ for every
$4\leq n\leq22$; for $17\leq n\leq22$ this follows by
block-diagonal containment from the $m=16$ padded cases.

\begin{proof}
First, $\mathcal P_{m,4}$ is an irreducible affine cone, as the
closure of the image of an irreducible affine space. It contains
an irreducible degree-$m$ polynomial: use the cyclic matrix with
diagonal entries $x_1$, entries $(i,i+1)$ equal to $x_2$ for
$i<m$, entry $(m,1)$ equal to $x_3$, and every other entry zero.
Its only perfect matchings are the diagonal and the full cycle,
so its permanent is $x_1^m+x_2^{m-1}x_3$.
Viewed in $\mathbb C[x_2,x_3][x_1]$, this polynomial is
Eisenstein at the prime $x_3$ and hence irreducible.
The reducible degree-$m$ forms constitute a closed subset:
projectivize each of the finitely many multiplication maps of
degrees $a$ and $m-a$, $1\leq a<m$, and take their closed images.
It follows that a general $q\in\mathcal P_{m,4}$ is irreducible.

Projectivizing the multiplication map
$(\ell,q)\mapsto\ell^{n-m}q$ gives a morphism
\[
 \mathbb P^3\times\mathbb P(\mathcal P_{m,4})
 \longrightarrow\mathbb P(\operatorname{Sym}^n\mathbb C^4)^*.
\]
Its image is closed because its domain is projective; its affine
cone is exactly the set in \eqref{R14:paddingvariety}, including
zero. For a generic irreducible $q$ of degree $m\geq2$, the
linear factor $\ell$ is uniquely determined by the product,
up to scalar, and then so is $q$. Thus the projective map is
generically injective. Its image dimension is
$3+(c_m-1)$ and the affine-cone dimension is $c_m+3$.

Every four-variable restriction of an invertible substitution
in $h_{m,n}$ has the form in \eqref{R14:paddingvariety}, so
closedness gives one inclusion of the projected orbit closure.
Conversely, for each permanent pencil $q=\Phi_m(A)$, any chosen
$\ell$ and its $m^2$ matrix-entry linear forms can be prescribed
simultaneously by an endomorphism substitution, since the
padding variable and the permanent variables are independent.
There is enough room because $m^2+1\leq n^2$.
Such substitutions lie in the orbit closure. Taking limits in
$q$ then shows that the whole closed set $\mathcal Z_{m,n}$ is
contained in the orbit closure, supported on its first four
variables. This proves equality of the projected image.

Let $\mathcal D_{n,4}$ be the determinant four-variable image
closure. By G-06, $\dim\mathcal D_{n,4}\leq2n^2+2$.
The coordinate rings of both cones are standard graded. The
Hilbert function of a nonzero standard graded affine cone of
dimension $c$ agrees for large $d$ with a polynomial of degree
$c-1$ and positive leading coefficient. Thus
\eqref{R14:criterion} implies
\[
 \dim\mathbb C[\mathcal Z_{m,n}]_d
 >\dim\mathbb C[\mathcal D_{n,4}]_d
\]
for every sufficiently large $d$.
Decompose these finite-dimensional $GL_4$ representations into
irreducibles. Their dimensions are the sums of multiplicities
times the positive integers $\dim S_\lambda\mathbb C^4$.
Hence some irreducible multiplicity on the first side is larger.

These $GL_4$ multiplicities agree with the corresponding
$GL_{n^2}$ multiplicities of length at most four in the original
orbit-closure coordinate rings. Indeed, highest-weight vectors
with those weights use only the first-four coefficient variables,
and the intersection of the full ideal with that polynomial
subring is the ideal of its projected image. The additional
raising operators involving later variables vanish on such
vectors. This proves the claimed multiplicity inequality.
G-05 states that the determinant multiplicity is already the
ambient maximum for weights of length at most three, so the
strict inequality must occur at length exactly four.

Substitution of $c_m=\binom{m+3}{3}$ into
\eqref{R14:criterion} gives $2n^2\leq\binom{m+3}{3}$, which is
\eqref{R14:explicitrange} and the displayed finite table.
Finally
\[
 \operatorname{per}\operatorname{diag}(L_m,zI_{n-m})
 =z^{n-m}\operatorname{per}(L_m)
\]
shows that the padded-permanent orbit closure is contained in
the size-$n$ permanent orbit closure. The resulting coordinate
quotient is surjective and therefore cannot increase
multiplicities. The padded gaps at $m=16$ thus imply the stated
unpadded gaps for $17\leq n\leq22$.
\end{proof}

If the expected formula
$c_m=\min\{\binom{m+3}{3},4m^2-2m+2\}$ were proved beyond the
certified range, \eqref{R14:criterion} would turn it into further
padded obstructions up to approximately $n=\sqrt2\,m$.
This expectation is a separate generic-fiber problem; neither
the finite certificates nor the dimension upper bound proves it.
The criterion above is unconditional whenever its stated
dimension lower bound has been established.

\subsection{R-15: two chords cross the determinant dimension for every size}

\paragraph{Statement R-15.}
For every $n\geq4$,
\begin{equation}\label{R15:strictdimension}
 \dim\mathcal P_{n,4}\geq2n^2+3
 >\dim\mathcal D_{n,4}.
\end{equation}
Consequently, for each fixed $n\geq4$ and every sufficiently
large coordinate-ring degree $d$, there is a partition
$\lambda\vdash nd$ with exactly four rows such that
\begin{equation}\label{R15:allngap}
 m_\lambda(\operatorname{per}_n)>m_\lambda(\det_n).
\end{equation}
The proof is algebraic and applies to all $n\geq4$ without an
extrapolation from the finite certificates. It supplies neither
an effective threshold for $d$ nor a specified $\lambda$ for
general $n$. In particular it is an unpadded, equal-size result,
and does not assert a growing-size determinantal-complexity lower
bound. At $n=4$, the lower bound is $35$, the full dimension of
the quaternary quartic space; it gives a purely algebraic proof
of the quartic dominance used in R-12.

\begin{proof}
Choose distinct complex numbers $a_1,\ldots,a_n$ and, separately,
distinct $b_1,\ldots,b_n$. Put
\[
 \ell_i=x-a_i y,\qquad v_i=z-b_i w,
 \qquad L=\prod_{i=1}^n\ell_i,
\]
with indices read cyclically modulo $n$.
Let $M_0$ have diagonal entries $\ell_i$, cycle entries
$(M_0)_{i,i+1}=v_i$, including $(M_0)_{n,1}=v_n$, and every
other entry zero. Write
$C_{ij}=\operatorname{per}(M_{0,\widehat i,\widehat j})$.
The image of the differential of the full permanent-pencil map
at $M_0$ is
\[
 W_0=\operatorname{span}\{xC_{ij},yC_{ij},zC_{ij},wC_{ij}:1\leq i,j\leq n\}.
\]
The complementary bipartite graph is a union of even paths,
each with a unique perfect matching. More explicitly, if $I$
is the cyclic interval from $j$ to $i-1$, of length
$k\in\{0,\ldots,n-1\}$, then
\begin{equation}\label{R15:cyclecofactor}
 C_{ij}=\prod_{h\in I}v_h
        \prod_{h\notin I\cup\{i\}}\ell_h.
\end{equation}

We first prove $\dim W_0\geq2n^2+2$.
Let $\Lambda_i\subset\mathbb P^3$ be the line
$\ell_i=v_i=0$, parametrized by
$(x,y,z,w)=(a_i u,u,b_i v,v)$.
Every cofactor from row $h\ne i$ has a factor $\ell_i$ or $v_i$
and restricts to zero on $\Lambda_i$.
The $n$ row-$i$ cofactors in \eqref{R15:cyclecofactor} restrict
to nonzero scalar multiples of
$u^{n-1-k}v^k$, $0\leq k<n$.
Multiplication by $y$ and $w$ therefore spans every degree-$n$
binary monomial on this line. Thus the restriction map
\[
 W_0\longrightarrow\bigoplus_{i=1}^n H^0(\Lambda_i,\mathcal O(n))
\]
is surjective, and its image has dimension $n(n+1)$.

For each proper nonempty cyclic interval $I$, together with
$I=\varnothing$ and $I=[n]$, put
\[
 Q_I=\prod_{i\in I}v_i\prod_{i\notin I}\ell_i.
\]
These polynomials belong to $W_0$: multiply the cofactor whose
omitted row immediately follows $I$ by the missing $\ell_i$;
for the two endpoint intervals use
$Q_\varnothing=\ell_iC_{ii}$ and
$Q_{[n]}=v_iC_{i,i+1}$.
Each $Q_I$ vanishes on every $\Lambda_i$.
They are linearly independent. For fixed length $k$ with
$1\leq k<n$, evaluate on the auxiliary line
$\Theta_t=\{\ell_t=v_{t+1}=0\}$.
The restriction of $Q_I$ is nonzero precisely when $I$ contains
$t$ and excludes $t+1$, which for a cyclic interval of that
length means that $I$ ends at $t$. The distinctness assumptions
ensure that the surviving scalar is nonzero. Thus these $n$
polynomials are independent for each $k$. Different lengths have
different bidegrees in $(x,y)$ and $(z,w)$, including the two
endpoint products. The kernel of the restriction map therefore
has dimension at least $n(n-1)+2$. Hence
\begin{equation}\label{R15:baserank}
 \dim W_0\geq n(n+1)+n(n-1)+2=2n^2+2.
\end{equation}

Next we identify enough of $W_0$ to exhibit a vector outside it.
The bidegree-$(n-2,2)$ component of $W_0$ is exactly
\begin{equation}\label{R15:edgespace}
 \left\{
 L\sum_{i=1}^n
 \frac{v_i g_i(z,w)}{\ell_i\ell_{i+1}}:
 g_i\in\operatorname{span}\{z,w\}
 \right\}.
\end{equation}
One inclusion follows by taking the cofactors whose interval has
length one and multiplying by $z,w$.
For the converse, a generator of this bidegree either has that
form, or comes from a length-two interval and multiplication
by a linear form $h(x,y)$.
Its rational expression after division by $L$ is
\[
 \frac{h(x,y)v_i v_{i+1}}{\ell_i\ell_{i+1}\ell_{i+2}}.
\]
Since $\ell_i$ and $\ell_{i+2}$ are independent, write
$h=A\ell_i+B\ell_{i+2}$ with scalars $A,B$.
The expression becomes the sum of an edge-$(i+1,i+2)$ term
with numerator $A v_i v_{i+1}$ and an edge-$(i,i+1)$ term
with numerator $B v_i v_{i+1}$. Both belong to
\eqref{R15:edgespace}. This proves equality; all generators of
the bidegree component have been included.

Define the degree-$n$ polynomial
\[
 T=\ell_1v_1v_3\prod_{i=4}^n\ell_i.
\]
We claim $T\notin W_0$. It has bidegree $(n-2,2)$, so
\eqref{R15:edgespace} is sufficient to test this claim.
On the affine chart $y=w=1$, put $t=x$, $u=z$.
Every rational expression in \eqref{R15:edgespace}, after
division by $L(t)$, has a partial fraction expansion
\[
 \sum_{i=1}^n\frac{R_i(u)}{t-a_i},
 \qquad R_i=E_i-E_{i-1},
\]
where $E_i(u)$ is a polynomial of degree at most two divisible
by $u-b_i$ and $E_0=E_n$. Indeed, take
$E_i=(u-b_i)g_i(u)/(a_i-a_{i+1})$ and use partial fractions
on each edge denominator.
For our particular polynomial,
\[
 \frac{T}{L}
 =\frac{(u-b_1)(u-b_3)}{(t-a_2)(t-a_3)}.
\]
Its residues vanish except at $a_2,a_3$, where they are
opposite nonzero scalar multiples of $(u-b_1)(u-b_3)$.
Thus $R_i=E_i-E_{i-1}$ would force all $E_i$ with $i\ne2$
to equal one common quadratic $E_*$. Divisibility at
$i=1,3,4$ gives three distinct roots $b_1,b_3,b_4$ of $E_*$,
so $E_*=0$.
Then $E_2$ is a nonzero scalar multiple of
$(u-b_1)(u-b_3)$, which is not divisible by $u-b_2$.
This contradiction proves $T\notin W_0$.

Now perturb $M_0$ by two entries:
\[
 (M_\varepsilon)_{1,3}=p=\varepsilon\ell_1,
 \qquad (M_\varepsilon)_{2,1}=q=\varepsilon v_3,
\]
leaving the other entries unchanged.
Let $C_{ij}(\varepsilon)$ be the complementary permanents at
this pencil and set $R=\prod_{i=4}^n\ell_i$.
Enumerating the two possible matchings in the indicated
complementary matrices gives
\[
 \begin{aligned}
 C_{31}(\varepsilon)&=(v_1v_2+p\ell_2)R,\\
 C_{32}(\varepsilon)&=(\ell_1v_2+pq)R,\\
 C_{33}(\varepsilon)&=(\ell_1\ell_2+v_1q)R.
 \end{aligned}
\]
Consequently the exact polynomial identity
\begin{equation}\label{R15:chordidentity}
 \ell_1C_{31}(\varepsilon)-v_1C_{32}(\varepsilon)
       -pC_{33}(\varepsilon)
       =-2\varepsilon^2T
\end{equation}
holds. Its left side belongs to the full Jacobian image
$W_\varepsilon$, since each multiplier is a linear form in
$x,y,z,w$. Thus $T\in W_\varepsilon$ whenever $\varepsilon\ne0$.

Choose $r=2n^2+2$ Jacobian columns independent at
$\varepsilon=0$, using \eqref{R15:baserank}.
Together with $T$ they are independent there, since
$T\notin W_0$. Some $(r+1)\times(r+1)$ coefficient minor of
these columns and $T$ is therefore nonzero at zero. It remains
nonzero for all $\varepsilon$ outside a finite set. For such
nonzero $\varepsilon$, $T$ already belongs to $W_\varepsilon$
by \eqref{R15:chordidentity}, so
$\dim W_\varepsilon\geq r+1=2n^2+3$.
The rank of a polynomial map's differential is a lower bound
for its image dimension, proving the first inequality in
\eqref{R15:strictdimension}. The second follows from G-06.

This construction can be made with integer coefficients.
For example take $a_i=i$, $b_i=i^2$.
Each Jacobian column has degree at most two in $\varepsilon$,
because only two matrix entries involve that parameter.
The augmented minor above has degree at most $2r=4n^2+4$ and
is not identically zero. Hence at least one integer
$\varepsilon\in\{1,\ldots,4n^2+5\}$ works.

Finally, the Hilbert-function and highest-weight argument in
R-14 applies to the two image cones in
\eqref{R15:strictdimension}. Their dimension inequality forces
a strict coordinate multiplicity inequality in every sufficiently
large degree. The determinant already has ambient multiplicities
for weights of length at most three by G-05, so the strict
inequality has length exactly four. This proves
\eqref{R15:allngap}.
\end{proof}

\subsection{R-16: simultaneous chords give an unbounded genuine-padding range}

\paragraph{Statement R-16.}
For $m\geq4$ put
\[
 h_m=\min\{\lfloor m/2\rfloor,m-3\},\qquad
 F_m=h_m(m-3).
\]
Thus $h_4=1$ and $h_m=\lfloor m/2\rfloor$ for $m\geq5$.
The quaternary permanent image satisfies the uniform bound
\begin{equation}\label{R16:dimension}
 \dim\mathcal P_{m,4}\geq 2m^2+2+F_m.
\end{equation}
At $m=4$ and $m=5$, this lower bound is respectively $35$ and
$56$, the whole ambient dimension. Thus it also proves the
quartic and quintic dominance of R-12 by an algebraic argument.
Consequently, if $n>m\geq4$ and
\begin{equation}\label{R16:paddingrange}
 2n^2\leq2m^2+F_m+2,
\end{equation}
then, in every sufficiently large coordinate-ring degree $d$,
some partition $\lambda\vdash nd$ with exactly four rows has
\[
 m_\lambda\bigl(\ell^{n-m}\operatorname{per}_m\bigr)
   >m_\lambda(\det_n).
\]
In particular, $n=m+1$ satisfies this conclusion for every
$m\geq12$. Together with the finite certificates in R-14,
this proves the conclusion for $n=m+1$ and every $m\geq8$.
For every fixed real $1<\rho<\sqrt5/2$, the same conclusion
holds for $n=\lfloor\rho m\rfloor$ and all sufficiently large
$m$. These are multiplicity gaps in a specified range of
linear padding; the proof does not specify the partition or
give a polynomial-time procedure to find it.

\begin{proof}
Fix a primitive $m$th root of unity $\zeta$ and use indices
$0,\ldots,m-1$ modulo $m$. Specialize the cyclic pencil of
R-15 to
\[
 \ell_i=x-\zeta^i y,\qquad v_i=z-\zeta^i w.
\]
Let $M_0$ have diagonal $\ell_i$ and cycle entries
$(M_0)_{i,i+1}=v_i$. Its differential image is denoted $W_0$.
The restriction-line and interval-product proof of R-15 applies
to this specialization, since both lists of roots are distinct,
and gives $\dim W_0\geq2m^2+2$.

We will simultaneously deform the disjoint row pairs
\[
 \mathcal A=\{0,2,\ldots,2(h_m-1)\},\qquad
 \{a,a+1\}\quad(a\in\mathcal A).
\]
For each $a\in\mathcal A$ add the two entries
\begin{equation}\label{R16:simultaneouspencil}
 (M_\varepsilon)_{a,a+2}=p_a=\varepsilon\ell_a,
 \qquad
 (M_\varepsilon)_{a+1,a}=q_a=\varepsilon v_{a+2}.
\end{equation}
All other entries remain those of $M_0$. The pairs of rows are
disjoint, including when the last triple of column indices
wraps around the cycle. Consequently rows $a,a+1$ have support
only in the three columns $a,a+1,a+2$, even when every other
pair has also been perturbed.

For $r\notin\{a,a+1\}$ let $R_{a,r}(\varepsilon)$ be the
permanent obtained by deleting rows $a,a+1,r$ and columns
$a,a+1,a+2$. Write $C_{r,j}(\varepsilon)$ for the complementary
permanent deleting row $r$ and column $j$. In each of the
three cases $j=a,a+1,a+2$, the two rows $a,a+1$ must match
the other two columns of that triple: they have no entries
outside it, and exactly two of its columns remain. The rest
of the matching is arbitrary in the same complementary
$(m-3)\times(m-3)$ matrix. Therefore, with no hypothesis on
the other row pairs,
\[
 \begin{aligned}
 C_{r,a}&=(v_a v_{a+1}+p_a\ell_{a+1})R_{a,r},\\
 C_{r,a+1}&=(\ell_a v_{a+1}+p_aq_a)R_{a,r},\\
 C_{r,a+2}&=(\ell_a\ell_{a+1}+v_aq_a)R_{a,r}.
 \end{aligned}
\]
Here and in the next display all permanents are evaluated at
$\varepsilon$. Subtraction gives the exact identity
\begin{equation}\label{R16:simultaneousidentity}
 \ell_a C_{r,a}-v_a C_{r,a+1}-p_a C_{r,a+2}
   =-2\varepsilon^2\ell_a v_a v_{a+2}R_{a,r}(\varepsilon).
\end{equation}
It follows that the polynomial vector
\[
 T_{a,k}(\varepsilon)
   =\ell_a v_a v_{a+2}R_{a,a+k}(\varepsilon)
\]
belongs to the full differential image $W_\varepsilon$ for
every $\varepsilon\ne0$ and $2\leq k\leq m-2$.
At $\varepsilon=0$ the remaining cycle paths have a unique
matching, so
\begin{equation}\label{R16:normalvectors}
 T_{a,k}:=T_{a,k}(0)
 =\ell_a v_a v_{a+2}
   \left(\prod_{i=a+2}^{a+k-1}v_i\right)
   \left(\prod_{i\notin[a,a+k]}\ell_i\right).
\end{equation}
Intervals and products are cyclic; the first product is empty
when $k=2$. This polynomial has bidegree $(m-k,k)$ in
$(x,y)$ and $(z,w)$.

We next prove that all $h_m(m-3)$ vectors in
\eqref{R16:normalvectors} are linearly independent modulo
$W_0$. This is the step that permits simultaneous rank growth,
rather than adding ranks from different deformations.

For $1\leq k<m$ set
\[
 F_{a,k}
   =\left(\prod_{i=a}^{a+k-2}v_i\right)
     \left(\prod_{i\notin[a,a+k-1]}\ell_i\right).
\]
The bidegree-$(m-k,k)$ part of $W_0$ is exactly
\begin{equation}\label{R16:bidegreespace}
 (W_0)_{m-k,k}
   =\operatorname{span}\{zF_{a,k},wF_{a,k}:0\leq a<m\}.
\end{equation}
Indeed the cycle cofactors with $k-1$ factors $v_i$ give
these generators directly. Every other generator of the
same bidegree is a cofactor with $k$ such factors multiplied
by a linear form $u(x,y)$. Write
$u=A\ell_a+B\ell_{a+k}$, possible because these two
linear forms are independent. The resulting polynomial is
$A v_aF_{a+1,k}+B v_{a+k-1}F_{a,k}$, proving the reverse
inclusion. This accounts for every cofactor and multiplier
of that bidegree. The two endpoint bidegrees have dimensions
at most $m+1$, so this argument also gives
$\dim W_0\leq2(m+1)+2m(m-1)=2m^2+2$.

Let $\tau$ act by
\[
 \tau(x,y,z,w)=(x,\zeta y,z,\zeta w).
\]
Then $F_{a,k}=\tau^aF_{0,k}$ and
$T_{a,k}=\tau^aT_{0,k}$. Equation
\eqref{R16:bidegreespace} shows that $W_0$ is invariant
under this cyclic action. Let $\pi_t$ denote its character
projection with eigenvalue $\zeta^t$.
We claim that
\begin{equation}\label{R16:allmodes}
 \pi_t(T_{0,k})\notin\pi_t(W_0)
 \quad\text{for every }2\leq k,t\leq m-2.
\end{equation}

Here is the full coefficient calculation proving the claim.
Write
\[
 \begin{aligned}
 A(x,y)&=\prod_{i=k}^{m-1}(x-\zeta^iy)
       =\sum_r A_r x^{m-k-r}y^r,\\
 B(z,w)&=\prod_{i=0}^{k-2}(z-\zeta^iw),\\
 E(z,w)&=\prod_{i=1}^{k-2}(z-\zeta^iw),\qquad B=(z-w)E.
 \end{aligned}
\]
Use subscripts for coefficients in the powers of the second
variable, with coefficients outside the degree range defined
to be zero. Every $A_r$, $0\leq r\leq m-k$, is nonzero.
For completeness the finite geometric-product identity gives
\[
 A_r=(-1)^r\zeta^{kr+r(r-1)/2}
      \prod_{j=1}^{r}
       \frac{1-\zeta^{m-k-r+j}}{1-\zeta^j}.
\]
This identity follows by induction on the number of linear
factors. Every exponent in numerator or denominator lies
strictly between $0$ and $m$, so no factor vanishes.

We have $F_{0,k}=AB$. The $(x,y)$ factor of $T_{0,k}$ is
$A(x,\zeta y)$, since its roots are $k+1,\ldots,m-1,0$.
If its $(z,w)$ factor is $B'(z,w)$, then
\[
 \begin{aligned}
 D(z,w):=B'(z,\zeta^{-1}w)
   &=(z-\zeta^{-1}w)(z-\zeta w)E(z,w)\\
   &=(z-w)B(z,w)+c\,zwE(z,w),\\
 c&=2-\zeta-\zeta^{-1}\ne0.
 \end{aligned}
\]
These equalities hold also for $k=2$, when $E=1$.

In character $t$, with $2\leq t\leq m-2$, the monomials
are precisely those with $r+s=t$; a second possibility
$r+s=t+m$ exceeds total degree $m$.
The allowed indices $s$ form the interval
\[
 \max\{0,t-(m-k)\}\leq s\leq\min\{k,t\}.
\]
It contains at least three consecutive integers, because
$2\leq k,t\leq m-2$.
The coefficient of such a monomial in $T_{0,k}$ is
$\zeta^t A_{t-s}D_s$. The two possible generators for
the character-$t$ part of \eqref{R16:bidegreespace}, namely
the projections of $zF_{0,k}$ and $wF_{0,k}$, have
coefficients $A_{t-s}B_s$ and $A_{t-s}B_{s-1}$.
Since all the displayed $A_{t-s}$ are nonzero, membership
in their span would imply that $D_s$ is a fixed linear
combination of $B_s$ and $B_{s-1}$ on the allowed interval.
Subtracting $(z-w)B$ and using $c\ne0$, it would follow
that $E_{s-1}$ is also such a linear combination.

That last assertion is impossible on any three consecutive
allowed indices $s=j,j+1,j+2$. To see this, use
$B_s=E_s-E_{s-1}$ to perform invertible column operations:
the relevant three-column determinant is nonzero if and
only if
\begin{equation}\label{R16:toeplitz}
 \det\bigl(E_{j+u-v}\bigr)_{0\leq u,v\leq2}\ne0.
\end{equation}
Here $0\leq j\leq k-2$. Put $N=k-2$.
Since $E_i=(-1)^i e_i(\zeta,\ldots,\zeta^N)$, row and
column signs change \eqref{R16:toeplitz} to the determinant
\[
 \det\bigl(e_{j+u-v}(\zeta,\ldots,\zeta^N)\bigr)_{0\leq u,v\leq2}
   =s_{(3^j)}(\zeta,\ldots,\zeta^N),
\]
by the dual Jacobi--Trudi identity (its representation-ring
form is equation~(2.4) in the attached proof of Sawin's Lemma~2.3;
R-B15). This Schur value is nonzero. If $N=0$, necessarily $j=0$ and the determinant
is one. Otherwise use the Weyl alternant at
$(1,\zeta,\ldots,\zeta^{N-1})$, then multiply by the
nonzero scalar $\zeta^{3j}$. For
$\lambda=(3^j,0^{N-j})$, the exponents
$\lambda_i+N-i$ are distinct integers in $[0,N+2]=[0,k]$.
Since $k<m$, their powers of $\zeta$ are distinct. The
alternant numerator is therefore a nonzero Vandermonde
determinant. Its denominator has the distinct exponents
$N-i$ and is nonzero as well. This proves
\eqref{R16:toeplitz}, and hence \eqref{R16:allmodes}.

Now fix $k$ and suppose
$\sum_{a\in\mathcal A}u_aT_{a,k}\in W_0$.
Projecting to each of the characters in
\eqref{R16:allmodes} gives
\[
 \sum_{a\in\mathcal A}u_a\zeta^{at}=0
 \qquad(2\leq t\leq m-2).
\]
There are $h_m\leq m-3$ distinct indices $a$.
Use the $h_m$ consecutive equations
$t=2,\ldots,h_m+1$. Their matrix is an ordinary
Vandermonde matrix in the distinct numbers $\zeta^a$,
with each column multiplied by the nonzero factor
$\zeta^{2a}$. It is invertible, so every $u_a=0$.
Thus the $h_m$ vectors for this $k$ are independent
modulo $W_0$. Different $k$ have different bidegrees
and $W_0$ is bigraded, so all $h_m(m-3)$ vectors are
independent modulo $W_0$.

Choose $2m^2+2$ differential columns at $M_0$ forming a
basis of $W_0$, and continue the same columns at
$M_\varepsilon$. Adjoin all the polynomial vectors
$T_{a,k}(\varepsilon)$. At zero these columns are
independent by the preceding argument, so one coefficient
minor of size $2m^2+2+F_m$ is a polynomial in
$\varepsilon$ with nonzero constant term. For generic
nonzero $\varepsilon$ it remains nonzero, and all its
adjoined vectors lie in $W_\varepsilon$ by
\eqref{R16:simultaneousidentity}. Hence
$\dim W_\varepsilon\geq2m^2+2+F_m$.
Differential rank bounds image dimension from below,
proving \eqref{R16:dimension}.

For padding, R-14 gives
$\dim\mathcal Z_{m,n}=\dim\mathcal P_{m,4}+3$.
Under \eqref{R16:paddingrange} this dimension is at least
$2n^2+3$, strictly larger than the determinant upper
bound $2n^2+2$. The Hilbert-function and highest-weight
argument of R-14 therefore supplies the asserted
exactly-four-row multiplicity gaps in every sufficiently
large degree. The explicit degree threshold of G-10
also applies to these pairs.

For $m\geq12$ one has
$F_m\geq(m-1)(m-3)/2\geq4m$, so
\eqref{R16:paddingrange} holds at $n=m+1$.
R-14 already covers each remaining $8\leq m\leq11$.
Finally, $F_m=m^2/2+O(m)$, and the right side of
\eqref{R16:paddingrange}, divided by $2m^2$, tends to
$5/4$. This proves the stated asymptotic range for every
fixed $\rho<\sqrt5/2$.
\end{proof}

\paragraph{What this does and does not identify.}
The simultaneous construction proves a lower bound on the
four-variable projection dimension, without assuming that
the row--column torus accounts for every generic fiber.
The generic maximal-rank formula
$\min\{\binom{m+3}{3},4m^2-2m+2\}$ is not used or proved.
The lower bound here is enough to retain four-row
multiplicity gaps under an unbounded range of genuine
padding, while the low-row exclusion of G-09 applies at
larger padding. No literature-priority claim is made for
these dimension deductions.


\subsection{R-17: exact determinant dimension and forced nonrectangular gaps}

\paragraph{Statement R-17.}
The four-variable determinant image has exact affine dimension
\begin{equation}\label{R17:detdimension}
 \dim\mathcal D_{n,4}=2n^2+2\qquad(n\geq1).
\end{equation}
For $n\geq3$ its affine $\mathrm{SL}_4$ quotient consequently
has dimension $2n^2-13$.
Moreover, at each of the padded pairs
\[
 (m,n)=(15,20),\qquad (m,n)=(16,22),
\]
there is a strictly positive multiplicity difference in some
\emph{nonrectangular} partition of exactly four rows in every
sufficiently large coordinate degree. In fact the effective
thresholds $d_{20}$ and $d_{22}$ of G-10 apply. At the second
pair, every positive-degree rectangular multiplicity on the
padded side is zero whenever the rectangular label is integral.

\begin{proof}
At the unperturbed cyclic pencil $M_0$ in R-15, each
complementary matrix has a unique perfect matching.
Therefore its determinant is its permanent multiplied by a
sign. The $(i,j)$ derivatives of the determinant additionally
have the fixed cofactor sign $(-1)^{i+j}$; hence each of the
four corresponding determinant-pencil Jacobian columns is
a sign times the permanent-pencil column. Their full images
are equal to $W_0$. The base-rank proof in R-15 gives
$\dim W_0\geq2n^2+2$ for $n\geq4$.
The determinant image dimension is at least this differential
rank and is at most $2n^2+2$ by G-06. This proves
\eqref{R17:detdimension} in that range. The remaining cases
are the full spaces of linear forms, quadrics, and cubic
surfaces, of dimensions $4,10,20$, respectively, as established
in R-10--R-11 and the elementary linear case.

The stable Fermat point and the $15$-dimensional quotient-fiber
argument of G-11 now give the asserted exact quotient dimension.
This is a direct consequence of the explicit cyclic pencil and
the established group-action bound; no claim of literature
priority for the dimension formula is made.

For the two padded pairs, R-13 and R-14 give
\[
 \begin{array}{c|c|c|c}
 (m,n)&\dim\mathcal P_{m,4}&\dim\mathcal Z_{m,n}
                    &\dim\mathcal D_{n,4}\\ \hline
 (15,20)&\binom{18}{3}=816&819&802\\
 (16,22)&\binom{19}{3}=969&972&970.
 \end{array}
\]
The strict dimension inequalities give exactly-four-row
multiplicity gaps by R-14, with the effective thresholds of
G-10. But $20=4\cdot15/3$, so the equality case of G-12
excludes a positive rectangular gap for the first pair.
For the second, $22>4\cdot16/3$, so the strict case of G-12
makes every positive-degree padded rectangular multiplicity
zero. If $4\nmid nd$, there is no integral four-row rectangle
of size $nd$ in the first place. Thus in every degree covered
by the Hilbert-function argument the required partition is
nonrectangular at both pairs.
\end{proof}


\subsection{R-18: two normal families improve the genuine-padding range}

\paragraph{Statement R-18.}
For $m\geq6$ put
\[
 j_m=\min\{\lfloor m/2\rfloor,m-5\},\qquad
 G_m=2j_m(m-5).
\]
Then
\begin{equation}\label{R18:dimension}
 \dim\mathcal P_{m,4}\geq2m^2+2+G_m.
\end{equation}
For any $n>m\geq6$ satisfying
\begin{equation}\label{R18:padding}
 2n^2\leq2m^2+G_m+2,
\end{equation}
R-14 and G-10 give a four-row multiplicity gap in every
coordinate degree $d\geq d_n$. In addition G-11 gives a
gap for the explicit rectangular partition
$\lambda=(nd/4)^4$ at every degree
$d=E_n k$, $k\geq K_n$, with its displayed effective
constants. In particular these conclusions hold eventually
for $n=\lfloor\rho m\rfloor$ and every fixed
$1<\rho<\sqrt{3/2}$.
The bounds of R-16 and R-18 are both valid; one may use the
larger lower bound at each particular size.

\begin{proof}
Keep the root-of-unity cyclic pencil and notation of R-16,
but take $j_m$ disjoint pairs with starting indices
$\mathcal A=\{0,2,\ldots,2(j_m-1)\}$ and use the entries
\[
 p_a=\varepsilon\ell_{a+1}\quad\hbox{at }(a,a+2),
 \qquad q_a=\varepsilon v_{a+2}\quad\hbox{at }(a+1,a).
\]
Each pair of rows still has support in just its three
columns. Therefore the simultaneous factorization in
R-16 remains exact. It gives two identities:
\begin{align*}
 \ell_a C_{r,a}-v_a C_{r,a+1}-p_a C_{r,a+2}
    &=-2p_aq_a v_a R_{a,r},\\
 q_a C_{r,a}+\ell_{a+1} C_{r,a+1}-v_{a+1} C_{r,a+2}
    &= 2p_aq_a\ell_{a+1}R_{a,r}.
\end{align*}
Both follow by substituting the three two-by-two permanent
factors already displayed in R-16; no other entries occur
in that substitution. Thus for $\varepsilon\ne0$ the
differential image contains the polynomial families
\[
 \begin{aligned}
 U_{a,k}(\varepsilon)
    &=\ell_{a+1}v_a v_{a+2}R_{a,a+k}(\varepsilon),\\
 V_{a,k}(\varepsilon)
    &=\ell_{a+1}^2v_{a+2}R_{a,a+k+1}(\varepsilon),
 \end{aligned}
\]
for $3\leq k\leq m-3$. Their values at zero, using the
same unique complementary matchings, are
\[
 \begin{aligned}
 U_{a,k}&=\ell_{a+1}v_a v_{a+2}
          \prod_{i=a+2}^{a+k-1}v_i
          \prod_{i\notin[a,a+k]}\ell_i,\\
 V_{a,k}&=\ell_{a+1}^2v_{a+2}
          \prod_{i=a+2}^{a+k}v_i
          \prod_{i\notin[a,a+k+1]}\ell_i.
 \end{aligned}
\]
They both have bidegree $(m-k,k)$, and their cyclic
translates are obtained by the same operator $\tau$ as
in R-16.

We prove that the two classes
$\pi_t(U_{0,k})$, $\pi_t(V_{0,k})$ are independent
modulo $\pi_t(W_0)$ whenever
\begin{equation}\label{R18:interior}
 3\leq k,t\leq m-3.
\end{equation}
This will give two independent Vandermonde systems.
Write $q=\zeta$, $K=q^k$, and $X=q^t$ just for the
coefficient calculation below. Retain
$A(x,y)=\prod_{i=k}^{m-1}(x-q^iy)$ from R-16, whose
coefficients $A_r$ are nonzero, and put
\[
 H(z,w)=\prod_{i=2}^{k-2}(z-q^iw),\qquad
 B(z,w)=(z-w)(z-qw)H(z,w).
\]
Thus $F_{0,k}=AB$. Define
\[
 \begin{aligned}
 \alpha&=\frac{K-q}{K-1},&
 \beta&=\frac{q-1}{K-1},\\
 \gamma_2&=\frac{q(K-1)}{qK-1},&
 \gamma_1&=\frac{(q-1)(K-q)}{(K-1)(qK-1)},&
 \gamma_0&=\frac{(q-1)^2}{(K-1)(qK-1)}.
 \end{aligned}
\]
All denominators are nonzero in the range
\eqref{R18:interior}.

The $(x,y)$ factors of $U_{0,k}$ and $V_{0,k}$ are
\[
 A_U=(x-qy)\prod_{i=k+1}^{m-1}(x-q^iy),\qquad
 A_V=(x-qy)^2\prod_{i=k+2}^{m-1}(x-q^iy).
\]
Their coefficient ratios are
\begin{equation}\label{R18:Aratios}
 \frac{(A_U)_r}{A_r}=\alpha q^r+\beta,
 \qquad
 \frac{(A_V)_r}{A_r}
    =\gamma_2 q^{2r}+\gamma_1 q^r+\gamma_0.
\end{equation}
Here is an explicit verification using the finite product
formula from R-16. If $u=q^r$, the first ratio is
$((K-q)u+q-1)/(K-1)$. For the second, put
$D=\prod_{i=k+2}^{m-1}(x-q^iy)$.
The same formula gives
\[
 \frac{D_r}{A_r}
  =\frac{(Ku-1)(qKu-1)}{(K-1)(qK-1)},\quad
 \frac{D_{r-1}}{A_r}
  =\frac{(u-1)(Ku-1)}{(K-1)(qK-1)},\quad
 \frac{D_{r-2}}{A_r}
  =\frac{(u-1)(u-q)}{q(K-1)(qK-1)}.
\]
Consequently $(A_V)_r/A_r$, obtained by subtracting
$2qD_{r-1}$ and adding $q^2D_{r-2}$, has numerator
\[
 (Ku-1)(qKu-1)-2q(u-1)(Ku-1)+q(u-1)(u-q).
\]
Its coefficients of $u^2,u,1$ are respectively
$q(K-1)^2$, $(q-1)(K-q)$, and $(q-1)^2$.
This proves \eqref{R18:Aratios}, including endpoint
coefficients interpreted as zero where necessary.

Let $v_b=z-q^bw$ in the next few displays. In character
$t$, after dividing every coordinate by the nonzero factor
$A_{t-s}$, the four vectors
$\pi_t(zF_{0,k})$, $\pi_t(wF_{0,k})$,
$\pi_t(U_{0,k})$, and $\pi_t(V_{0,k})$ have the coefficient
sequences of $H$ times the following four binary cubics:
\begin{equation}\label{R18:cubics}
 zv_0v_1,\quad wv_0v_1,\quad
 \mathcal U_X=\alpha Xv_{-1}v_1^2+\beta v_0v_2v_{k-1},\quad
 \mathcal V_X=\gamma_2X^2v_0^2v_1
               +\gamma_1Xv_1^2v_{k-1}
               +\gamma_0v_2v_{k-1}v_k.
\end{equation}
To check the shifts in this assertion, the $(z,w)$ factor
of $U_{0,k}$, after replacing $w$ by $q^{-1}w$, is
$v_{-1}v_1^2H$. That factor with $w$ replaced back by
$qw$ becomes $v_0v_2v_{k-1}H$.
The $(z,w)$ factor of $V_{0,k}$, after replacing $w$ by
$q^{-2}w$, is $v_0^2v_1H$; its next two shifts are
$v_1^2v_{k-1}H$ and $v_2v_{k-1}v_kH$.
Together with $q^r=q^{t-s}$ and
\eqref{R18:Aratios}, these identities give exactly
\eqref{R18:cubics}.

These four cubics are linearly independent. Modulo the
first two, it suffices to evaluate at $z/w=1,q$, because
the kernel of those two evaluations is precisely
$v_0v_1\operatorname{span}\{z,w\}$.
Direct substitution and simplification give
\begin{equation}\label{R18:evaldet}
 \begin{vmatrix}
 \mathcal U_X(1,1)&\mathcal V_X(1,1)\\
 \mathcal U_X(q,1)&\mathcal V_X(q,1)
 \end{vmatrix}
 =\frac{(q-K)(q^2-K)(q-1)^6(q+1)}
        {q(K-1)(qK-1)}\ne0.
\end{equation}
The cancellation of $X$ in this expression is exact.
Nonvanishing follows from the order $m$ of $q$ and
$3\leq k\leq m-3$; in particular $q\ne\pm1$ and
$K\notin\{1,q,q^2,q^{-1}\}$.

Under \eqref{R18:interior}, the allowed coefficient
indices $s$ contain at least four consecutive values
$b,b+1,b+2,b+3$, with $0\leq b\leq k-3$.
Multiplication by $H$ takes the cubic monomial basis
to four columns whose restriction to those coordinates
has determinant
\[
 \det(H_{b+u-v})_{0\leq u,v\leq3}.
\]
Up to row and column signs this is the dual
Jacobi--Trudi determinant
\[
 s_{(4^b)}(q^2,q^3,\ldots,q^{k-2}).
\]
It is nonzero by the same alternant argument as in
R-16: there are $N=k-3$ variables and the numerator
exponents for $(4^b)$ are distinct in
$[0,N+3]=[0,k]$, strictly shorter than the order $m$.
For $N=0$, necessarily $b=0$ and the determinant is one.
Thus restriction to the four selected coefficients
preserves independence of the cubics
\eqref{R18:cubics}. This proves the asserted independence
of the two normal classes in every character
\eqref{R18:interior}.

Now suppose, for fixed $k$, that a linear combination
\[
 \sum_{a\in\mathcal A}
        (c_aU_{a,k}+d_aV_{a,k})
\]
belongs to $W_0$. Projecting to character $t$ gives a
linear combination of two independent normal classes,
so both $\sum_a c_aq^{at}$ and
$\sum_a d_aq^{at}$ must vanish.
Since $j_m\leq m-5$, the $j_m$ consecutive values
$t=3,\ldots,j_m+2$ are available. Each of these two
systems is an invertible Vandermonde system in the
distinct $q^a$. Hence every coefficient vanishes.
The $2j_m$ normal vectors are independent for this
bidegree. Different $k$ have different bidegrees, so
the full collection has $2j_m(m-5)=G_m$ independent
classes modulo $W_0$.

Adjoin their polynomial deformations to a basis of
the $2m^2+2$ cyclic Jacobian columns. At zero the
augmented columns are independent. A nonzero coefficient
minor remains nonzero for generic $\varepsilon$, when
all these columns lie in the actual deformed Jacobian
image by the two exact cofactor identities. This proves
\eqref{R18:dimension} by differential rank.

The padding conclusion follows from R-14 and G-10.
For the rectangular conclusion, note that
$G_m\leq m(m-5)$, so \eqref{R18:padding} implies
$n^2<3m^2/2$ and hence $n<4m/3$.
Both hypotheses of G-11 therefore hold.
Finally $G_m=m^2+O(m)$, so the upper admissible ratio
in \eqref{R18:padding} tends to $\sqrt{3/2}$.
\end{proof}


\subsection{R-19: simultaneous row and column blocks improve the range again}

\paragraph{Statement R-19.}
For $m\geq10$ put
\[
 a_m=\min\{\lfloor m/3\rfloor,m-9\},\qquad
 H_m=4a_m(m-9).
\]
Then
\begin{equation}\label{R19:dimension}
 \dim\mathcal P_{m,4}\geq2m^2+2+H_m.
\end{equation}
If $n>m\geq10$ and
\begin{equation}\label{R19:padding}
 2n^2\leq2m^2+H_m+2,
\end{equation}
the four-row multiplicity gaps of R-14/G-10 and the explicit
rectangular gaps of G-11 both hold. In particular, they hold
eventually for $n=\lfloor\rho m\rfloor$ and every fixed
$1<\rho<\sqrt{5/3}$. Earlier lower bounds remain available
when they are larger at a particular finite size.

\begin{proof}
Use the same cyclic pencil with
$\ell_i=x-\zeta^iy$, $v_i=z-\zeta^iw$.
Choose centers
\[
 \mathcal A=\{1,4,\ldots,1+3(a_m-1)\}.
\]
The triples of rows $a-1,a,a+1$ are disjoint and any two
centers have cyclic distance at least three.
For each center add
\begin{equation}\label{R19:threechords}
 p_a=\varepsilon\ell_{a+1}\text{ at }(a,a+2),\qquad
 q_a=\varepsilon v_{a+2}\text{ at }(a+1,a),\qquad
 r_a=\varepsilon\ell_a\text{ at }(a-1,a+1).
\end{equation}
The two rows $a,a+1$ still have support only in columns
$a,a+1,a+2$. The two columns $a,a+1$ have support only
in rows $a-1,a,a+1$. To verify the latter assertion for
all simultaneous motifs, an additional forward chord
from another center $b$ could enter these columns only
if $b=a-2$ or $a-1$, and an additional backward or third
chord only if $b=a-1$ or $a+1$. These cases are excluded
by the spacing. The analogous row-support assertion is
immediate from the disjoint row triples.

Thus the two row identities of R-18 remain exact and
give the same families $U_{a,k}(\varepsilon)$ and
$V_{a,k}(\varepsilon)$, with the remaining complementary
permanents now evaluated at the three-chord pencil.
There are also two exact column identities. For a column
$c\notin\{a,a+1\}$ let $S_{a,c}(\varepsilon)$ be the
permanent obtained by deleting rows $a-1,a,a+1$ and
columns $a,a+1,c$. The complementary permanents factor as
\[
 \begin{aligned}
 C_{a-1,c}&=(\ell_a\ell_{a+1}+v_aq_a)S_{a,c},\\
 C_{a,c}&=(v_{a-1}\ell_{a+1}+q_ar_a)S_{a,c},\\
 C_{a+1,c}&=(v_{a-1}v_a+\ell_ar_a)S_{a,c}.
 \end{aligned}
\]
Indeed the two remaining rows of the triple must match
the two supported columns; all other matching choices
occur in the same complementary matrix.
Consequently
\begin{align*}
 v_{a-1}C_{a-1,c}-\ell_a C_{a,c}-q_a C_{a+1,c}
    &=-2q_ar_a\ell_a S_{a,c},\\
 r_a C_{a-1,c}+v_a C_{a,c}-\ell_{a+1}C_{a+1,c}
    &= 2q_ar_av_a S_{a,c}.
\end{align*}
After division by the nonzero scalar $\varepsilon^2$,
these show that the differential image contains
\[
 C_{a,k}(\varepsilon)=\ell_a^2v_{a+2}S_{a,a-k}(\varepsilon),
 \qquad
 D_{a,k}(\varepsilon)=\ell_av_av_{a+2}S_{a,a-k+1}(\varepsilon).
\]
We use $5\leq k\leq m-5$ below.
At zero the unique cycle matching gives
\[
 \begin{aligned}
 C_{a,k}&=\ell_a^2v_{a+2}
       \prod_{i=a-k}^{a-2}v_i
       \prod_{i\notin[a-k,a+1]}\ell_i,\\
 D_{a,k}&=\ell_av_av_{a+2}
       \prod_{i=a-k+1}^{a-2}v_i
       \prod_{i\notin[a-k+1,a+1]}\ell_i.
 \end{aligned}
\]
All four families $U,V,C,D$ have bidegree $(m-k,k)$,
and translating their center by $a$ applies $\tau^a$.

We prove fourfold normal independence in each character
\begin{equation}\label{R19:interior}
 5\leq k,t\leq m-5.
\end{equation}
As in R-18 write $q=\zeta$, $K=q^k$, $X=q^t$, and
divide character-$t$ coordinates by the nonzero
coefficients $A_{t-s}$ of
$A(x,y)=\prod_{i=k}^{m-1}(x-q^iy)$.
Set
\[
 J(z,w)=\prod_{i=2}^{k-4}(z-q^iw),\qquad
 v_b=z-q^bw.
\]
The two base columns and the two row-family columns,
after this normalization, are $J$ times
$v_{k-3}v_{k-2}$ times the four independent cubics
of \eqref{R18:cubics}. Thus these four quintics span
\begin{equation}\label{R19:rowquintics}
 v_{k-3}v_{k-2}\operatorname{Sym}^3\mathbb C^2.
\end{equation}

We compute the two column-family sequences explicitly.
Put $M=m-k$ and
$\widetilde A=A(x,q^{2-k}y)=\prod_{i=2}^{M+1}(x-q^iy)$.
Their $(x,y)$ factors at center zero are
\[
 A_C=(x-y)^2\prod_{i=2}^{M-1}(x-q^iy),\qquad
 A_D=(x-y)\prod_{i=2}^{M}(x-q^iy).
\]
Define
\[
 \begin{aligned}
 \delta_0&=\frac{K(q-1)}{q(K-1)},&
 \delta_1&=\frac{K-q}{q(K-1)},\\
 \eta_0&=\frac{K^2(q-1)^2}{q(K-1)(qK-1)},&
 \eta_1&=\frac{K(K-q)(q-1)}{q(K-1)(qK-1)},&
 \eta_2&=\frac{K-1}{qK-1}.
 \end{aligned}
\]
The following are polynomial identities:
\[
 \begin{aligned}
 A_D&=\delta_0\widetilde A(x,y)
           +\delta_1\widetilde A(x,q^{-1}y),\\
 A_C&=\eta_0\widetilde A(x,y)
           +\eta_1\widetilde A(x,q^{-1}y)
           +\eta_2\widetilde A(x,q^{-2}y).
 \end{aligned}
\]
For the first, cancel the common factors with indices
$2,\ldots,M$ and compare the remaining linear forms.
For the second cancel those with indices $2,\ldots,M-1$;
the remaining identity is
\[
 \eta_0(x-K^{-1}y)(x-qK^{-1}y)
 +\eta_1(x-qy)(x-K^{-1}y)
 +\eta_2(x-y)(x-qy)=(x-y)^2,
\]
which follows by its three coefficients. This verifies
the coefficient ratios without any genericity assumption.

Since $\widetilde A_r=q^{(2-k)r}A_r$,
the two column vectors, after the common coordinate
division and multiplication by the nonzero character
scalar $X^{k-2}$, are $J$ times the quintics
\begin{align*}
 \mathcal C_X={}&
 \eta_0v_{-2}v_{-1}v_0v_1v_k
 +\eta_1X^{-1}v_{-1}v_0v_1v_{k-3}v_{k+1}
 +\eta_2X^{-2}v_0v_1v_{k-3}v_{k-2}v_{k+2},\\
 \mathcal D_X={}&
 \delta_0v_{-1}v_0v_1v_{k-2}v_k
 +\delta_1X^{-1}v_0v_1v_{k-3}v_{k-1}v_{k+1}.
\end{align*}
For example the $(z,w)$ factor of $C_{0,k}$ has roots
$2,m-k,\ldots,m-2$; shifting $w$ by $q^{k-2}$ changes
them to $k,-2,-1,0,\ldots,k-4$.
Its next two shifts give the other two terms above.
The factor of $D_{0,k}$ has roots
$0,2,m-k+1,\ldots,m-2$; its corresponding two shifts
give the terms of $\mathcal D_X$. This checks every
factor and character scalar in these expressions.

To test independence from \eqref{R19:rowquintics},
evaluate the two new quintics at
$z/w=K/q^3$ and $K/q^2$, the distinct roots of
$v_{k-3}v_{k-2}$. Substitution gives
\begin{equation}\label{R19:evaldet}
\begin{aligned}
 &\begin{vmatrix}
 \mathcal C_X(K/q^3,1)&\mathcal D_X(K/q^3,1)\\
 \mathcal C_X(K/q^2,1)&\mathcal D_X(K/q^2,1)
 \end{vmatrix}\\
 &\qquad=
 \frac{K^6(q-K)(q^2-K)^2(q^3-K)^2(q^4-K)
             (q-1)^6(q+1)(q^2+q+1)}
      {q^{27}(K-1)(qK-1)}\ne0.
\end{aligned}
\end{equation}
Here too the cancellation of $X$ is exact.
Every displayed factor is nonzero: $m\geq10$ and
$5\leq k\leq m-5$, so $q$ has order neither two nor
three and $K$ differs from $1,q,q^2,q^3,q^4,q^{-1}$.
The four quintics in \eqref{R19:rowquintics} and the
two new ones consequently form a basis of all binary
quintics.

Under \eqref{R19:interior}, the allowed indices $s$
contain six consecutive values $b,\ldots,b+5$, where
$0\leq b\leq k-5$. Multiplication by $J$ and restriction
to these six coordinates has, in the quintic monomial
basis, determinant
\[
 \det(J_{b+u-v})_{0\leq u,v\leq5}
   =\pm s_{(6^b)}(q^2,q^3,\ldots,q^{k-4})\ne0.
\]
There are $N=k-5$ variables. The Weyl numerator
exponents for $(6^b)$ are distinct in
$[0,N+5]=[0,k]$, below the order $m$, so the same
Vandermonde proof as in R-16 applies. When $N=0$,
the determinant is one. Thus all six normalized
character columns are independent; equivalently,
the four classes of $U,V,C,D$ are independent modulo
the base character space.

A relation among their translates at centers in
$\mathcal A$ now gives four separate equations
$\sum_{a\in\mathcal A}c_aq^{at}=0$ in each of the
available characters. Since $a_m\leq m-9$, the
$a_m$ consecutive characters
$t=5,\ldots,a_m+4$ lie in the allowed range.
Each system is an invertible Vandermonde system in
the distinct $q^a$. Hence the $4a_m$ classes are
independent in each bidegree. Different bidegrees
are independent, giving $4a_m(m-9)=H_m$ classes
modulo $W_0$.

All four families have polynomial deformations lying
in the actual differential image for nonzero
$\varepsilon$. The augmented-minor argument of
R-16, with the $2m^2+2$ base columns, proves
\eqref{R19:dimension}.
R-14 and G-10 give the general four-row consequence.
Furthermore $H_m\leq(4m/3)(m-9)$, so
\eqref{R19:padding} implies
$n^2\leq5m^2/3-6m+1<16m^2/9$.
Thus $n<4m/3$ and G-11 also applies.
Finally $H_m=4m^2/3+O(m)$, giving the asserted
asymptotic ratio $\sqrt{5/3}$.
\end{proof}


\subsection{What the attached GCT IV, VII, and VIII positivity claims supply}

The attached GCT~IV is the later Blasiak--Mulmuley--Sohoni manuscript
corresponding to their 2015 memoir, with filename
\texttt{Apr13KroneckerGCT4.tex}. Its \S12.1 theorem
(\texttt{t nonstandard schur-weyl duality}) proves nonstandard
Schur--Weyl duality for general $\dim V,\dim W$; the later
explicit-basis limitation should not obscure this proved
structural result. Its Theorem~15.21(vii)--(viii)
(\texttt{t main canonical basis}) proves a highest-weight crystal
count for $g(\alpha,\beta,\gamma)$ when
$\ell(\alpha),\ell(\beta)\leq2$, including the symmetric/exterior
split. Theorem~18.2 (\texttt{t consequence for Kronecker}) makes
the count explicit, and Corollary~18.4 gives the generating
function used in R-10. This supplies a positive rule for rectangles
of height two, or width two after simultaneous conjugation.
The construction for the full tensor power remains Conjecture~19.1
(\texttt{cj canonical basis X\^{}r}); \S19.2 describes the failed
flatness of a naive extension beyond the two-row case.

Even in the proved case, counting highest-weight elements does not
mean that every matrix coefficient is nonnegative. GCT~IV,
Example~17.8 (\texttt{ex not positive}), displays the divided-power
coefficient $q^{10}-q^4-q^{-4}+q^{-10}$. This concerns the specified
basis and operator; it is a concrete warning against replacing a
precise canonical-basis positivity statement by unrestricted
coefficient positivity.

GCT~VII (Mulmuley, 2008a, revised extended abstract) proves in
Theorem~1.2 (\texttt{tnonstd}) existence of the nonstandard
quantization, complete reducibility of its polynomial
representations, and its Peter--Weyl statement. The nonstandard
duality and reciprocity needed to recover plethysm are explicitly
Conjectures~1.3--1.4 (\texttt{cduality}, \texttt{creciprocity});
the refined form is Conjecture~6.1 (\texttt{crecirefined}).
The proposed formula
$a_{\lambda,\mu}^{\pi}=\sum_\alpha
m_\lambda^\alpha n_\pi^\alpha$ is conditional on those structures.
The introduction states that correctness of the basis algorithms
and their required properties has not been proved there.

GCT~VIII (Mulmuley, 2008b, revised extended abstract) makes the
additional gates explicit: local crystal existence is Conjecture~2.2
(\texttt{clocalcrystalb}); minimum-degree globalization is
Conjectures~2.5, 3.1--3.2; cellular decomposition is Conjectures~2.6,
3.3 and the specialization-compatible quasi-cell condition is
Conjecture~3.7. Kronecker positivity is Conjectures~2.7 and~3.4,
with specified signs and powers of $q-q^{-1}$ allowed. The general
plethysm setting uses the weaker saturation and nonstandard
positivity Conjectures~2.8--2.9 and~3.5--3.6. Finally, \S6
(\texttt{scomplexity}) also requires efficient combinatorial
labels and crystal operations to obtain a $\#P$ formula; these
complexity properties are themselves proposed rather than proved.

Its discussion of the Riemann hypothesis in the introduction and
\S2.4.2 is a suggested extension of finite-field purity and
perverse-sheaf machinery to the nonflat, noncommutative objects at
issue. It is not a theorem deriving these conjectures from the
classical RH or Dirichlet GRH. Accepting the stipulated analytic
theorems therefore leaves specific algebraic gates: construct the
required realization, prove its cellular and specialization
properties, and control the sizes of its labels and operations.
The supplied papers prove the height-two case and propose this
general route; they supply no unconditional constructive positive
formula for arbitrary rectangles and no theorem controlling the
orbit-closure ranks in R-05--R-06.

\paragraph{Algebraic-geometric dictionary for R-01--R-08.}
R-01 and R-02 are identities of Schur operations and restriction maps
on representation $K_0$, evaluated over $\mathbb C$ for Haar integration.
The tensor homomorphism is the integral group-scheme morphism
$GL_a\times GL_n\to GL_{an}$, and its kernel contains the central
$\mathbb G_m$ described in R-03. R-04 is the cycle-index expression
for symmetric powers in a $\lambda$-ring; the Adams operations preserve
the representation ring, not its cone of actual representations.
For R-05--R-06 the intrinsic algebraic object is the graded orbit-map
image, or the coefficient map in \eqref{R06:coefficients}. Its rank
strata are defined by determinantal (Fitting) ideals. The Gaussian
and Hermitian metrics are analytic devices for the same rank; they
are not asserted to be motives. The Schur idempotent is defined over
$\mathbb Z[1/D!]$, and the integral matrix in \eqref{R06:integer}
avoids treating characteristic-zero semisimplicity as an integral
direct-sum theorem. R-07 is a statement about restrictions of locally
free Schur constructions to the binary-form subscheme and the
scheme-theoretic image of the linear-substitution morphism. Its
formulas can be cleared of their binomial denominators over
$\mathbb Z$, while the basis and orbit-closure claims here are in
characteristic zero. R-08 is additivity in equivariant $K_0$ applied
to a resolution; the numerator remembers the quotient's equations.
Exact SGA/EGA/FGA and cote~133 locators for this dictionary are supplied
in the separate Grothendieck part of the paper.
R-09 identifies a four-factor invariant of
$\operatorname{Sym}(V_1\otimes V_2\otimes V_3\otimes V_4)$ with the
composition of two tensor-product restriction functors; its degree
$nd$ determinant-isotypic coefficient is the fourfold Kronecker
multiplicity. The finite CUE cutoff is the rank restriction on a
Schur functor, not an arithmetic purity hypothesis.

\subsection{R-20: a five-row certificate reaches the rectangular padding boundary}

\paragraph{Statement R-20.}
Put $M_0=8{,}000{,}001$. For every integer $m\geq M_0$, let
\[
 h_m=\left\lfloor\frac{m-97}{5}\right\rfloor,
 \qquad J_m=8h_m(m-317).
\]
Then
\begin{equation}\label{R20:dimension}
 \dim\mathcal P_{m,4}\geq2m^2+2+J_m.
\end{equation}
Consequently, for every pair of integers
\begin{equation}\label{R20:rectrange}
 m\geq M_0,\qquad m<n<4m/3,
\end{equation}
the explicit rectangular multiplicity gaps of G-11 hold, with its
labels $\lambda=(nd/4)^4$, degrees $d=E_nk$, and $k\geq K_n$.
Together with G-12, this makes $n<4m/3$ the exact eventual range,
among $n>m$, in which a rectangular multiplicity gap exists.
This is a statement about the specified four-row sector and its
large coordinate degrees.

More generally, whenever $m\geq M_0$, $n>m$, and
\begin{equation}\label{R20:allrange}
 2n^2\leq2m^2+J_m+2,
\end{equation}
R-14/G-10 give an exactly-four-row multiplicity gap in every degree
$d\geq d_n$. If also $n\geq4m/3$, every such gap is nonrectangular
by G-12. In particular, for every fixed $4/3<\rho<3/\sqrt5$, it includes
$n=\lfloor\rho m\rfloor$ for all sufficiently large $m$.
No maximal generic-rank formula is assumed.

\begin{proof}
We prove the dimension bound by an explicit differential certificate.

\paragraph{1. One simultaneous local construction.}
Let $q$ be a primitive $m$th root, to be chosen more carefully below,
and put
\[
 \ell_i=x-q^iy,\quad v_i=z-q^iw,\qquad
 L_c=x-cy,\quad V_c=z-cw.
\]
Start with the cyclic pencil of R-15. At a block starting at zero,
replace the five rows $0,\ldots,4$, in columns $0,\ldots,5$, by
\begin{equation}\label{R20:motif}
 \begin{pmatrix}
 \ell_0&v_0&\varepsilon L_2&\varepsilon L_3&\varepsilon V_3&0\\
 0&\ell_1&v_1&\varepsilon L_4&\varepsilon L_3&\varepsilon V_2\\
 0&0&\ell_2&v_2&\varepsilon V_6&0\\
 0&\varepsilon V_5&0&\ell_3&v_3&\varepsilon V_4\\
 0&\varepsilon L_2&\varepsilon V_7&\varepsilon L_5&\ell_4&v_4
 \end{pmatrix}.
\end{equation}
At starting positions $a=0,5,\ldots,5(h_m-1)$ use the translated
motif: shift all row and column indices by $a$ and replace
$(y,w)$ by $(q^ay,q^aw)$ in its entries. The five-row blocks are
disjoint. The four rows $a+1,\ldots,a+4$ have support only in
columns $a+1,\ldots,a+5$. The four columns
$a+1,\ldots,a+4$ have support only in rows $a,\ldots,a+4$.
Both assertions hold with every block present, since a neighboring
block shares only the boundary column. All remaining entries are
those of the original cycle.

Here is the local cofactor calculation used for both supports.
For a $4\times5$ matrix $B(\varepsilon)$ with diagonal entries
$a_i$, superdiagonal entries $b_i$, and other entries
$\varepsilon c_{ij}$, form $\widetilde B$ by multiplying the
off-bidiagonal entry $(i,j)$ by
\[
 \sigma_{ij}=\begin{cases}
 (-1)^{j-i-1},&j>i+1,\\
 (-1)^{i-j},&j<i.
 \end{cases}
\]
Write $P_j=\operatorname{per}(B\text{ with column }j\text{ deleted})$
and $D_j=\det(\widetilde B\text{ with column }j\text{ deleted})$.
There is one matching at order zero. A matching with one additional
edge $(i,j)$ has permutation sign $(-1)^{j-i-1}$ for $j>i+1$,
or $(-1)^{i-j}$ for $j<i$: the forced remaining bidiagonal edges
form the corresponding consecutive cycle. Thus
\begin{equation}\label{R20:secondorder}
 P_j-D_j=\varepsilon^2 E_j+O(\varepsilon^3),
 \qquad E_0=E_4=0.
\end{equation}
For completeness, the entire second-order expansion is
\begin{align*}
 E_1/2={}&c_{02}c_{10}b_2b_3
 +c_{03}(c_{10}a_2b_3+c_{20}b_1b_3)\\
 &+c_{04}(c_{10}a_2a_3+c_{20}b_1a_3+c_{30}b_1b_2),\\
 E_2/2={}&c_{13}(c_{21}a_0b_3+c_{20}b_0b_3)\\
 &+c_{14}(c_{21}a_0a_3+c_{31}a_0b_2
                   +c_{20}b_0a_3+c_{30}b_0b_2)\\
 &+c_{03}c_{20}a_1b_3
 +c_{04}(c_{20}a_1a_3+c_{30}a_1b_2),\\
 E_3/2={}&c_{24}(c_{32}a_0a_1+c_{31}a_0b_1+c_{30}b_0b_1)\\
 &+c_{14}(c_{31}a_0a_2+c_{30}b_0a_2)
 +c_{04}c_{30}a_1a_2.
\end{align*}
These identities can be verified directly by the $24$ permutations
of each $4\times4$ minor: the displayed $6,9,6$ terms are precisely
the two-additional-edge matchings whose signed determinant term is
negative. This specifies the finite calculation completely.

The determinant row relation gives
\begin{equation}\label{R20:relation}
 \sum_{j=0}^4(-1)^{i+j}\widetilde B_{ij}P_j
 =\varepsilon^2(a_iE_i-b_iE_{i+1})+O(\varepsilon^3).
\end{equation}
Apply it first to rows $1,\ldots,4$ and columns $1,\ldots,5$
of \eqref{R20:motif}, and then to the transpose of its submatrix
in rows $0,\ldots,4$ and columns $1,\ldots,4$.
In the global complementary permanents these identities have a
common remaining factor. For an external deleted row $s$, the
first set factors because the four supported rows must occupy the
four remaining supported columns. For an external deleted column
$c$, the second set factors because the four supported columns
must occupy the four remaining supported rows. The presence of
other blocks does not change either argument.

The coefficient of $\varepsilon^2$, divided by two, has the following
highest components in $(z,w)$ degree. This list follows by substituting
\eqref{R20:motif} into the preceding complete expansion:
\begin{align*}
 R_0={}&-v_1v_2(L_3V_5v_4+V_2V_5\ell_4+V_2L_2v_3),\\
 R_1={}&-v_1v_2v_4V_5V_6,\\
 R_2={}&\ell_3v_1v_4V_5V_6
          -v_2v_3V_4(\ell_1V_7+L_2v_1),\\
 R_3={}&\ell_4v_2V_4(\ell_1V_7+L_2v_1),\\[2pt]
 S_0={}&-\ell_1L_2v_2(\ell_4V_5+L_2v_3),\\
 S_1={}&v_1L_2v_2(\ell_4V_5+L_2v_3)
          -\ell_2L_4v_0v_3V_7,\\
 S_2={}&L_4v_0v_2v_3V_7,\\
 S_3={}&v_3(L_5V_6v_0v_1+L_3V_7v_0v_2
                +\ell_1V_7V_3v_2+L_2V_3v_1v_2).
\end{align*}
Their $(z,w)$ degrees are respectively
\begin{equation}\label{R20:grades}
 (4,5,4,3)\quad\hbox{and}\quad(2,3,4,4).
\end{equation}
Every listed polynomial has total degree five. The complete
coefficient in \eqref{R20:relation} may have additional components
of strictly smaller $(z,w)$ degree; those will be handled below,
without projecting an actual differential image by an unjustified
operation.

At $\varepsilon=0$ the common outside factors are exactly
\begin{equation}\label{R20:outside}
 T_s^{\rm row}=\ell_0\prod_{i=5}^{s-1}v_i
                         \prod_{i=s+1}^{m-1}\ell_i,
 \qquad
 T_c^{\rm col}=\prod_{i=5}^{c-1}\ell_i
                         \prod_{i=c}^{m-1}v_i.
\end{equation}
This follows either by the unique cycle matching or by successively
matching the forced boundary column. If a local polynomial has
$(z,w)$ degree $g$, choose
\begin{equation}\label{R20:alignment}
 s=k+5-g\quad\hbox{for a row family},\qquad
 c=m+g-k\quad\hbox{for a column family}.
\end{equation}
Its highest component times \eqref{R20:outside} then has bidegree
$(m-k,k)$. All deleted indices are external and distinct as required
when $9\leq k\leq m-9$.

\paragraph{2. Fixed rational-function matrices controlling all modes.}
Write $K=q^k$, $X=q^t$. The following construction defines two
$6\times6$ matrices over $\mathbb Q(q,K,X)$, using only the eight
displayed degree-five polynomials. It therefore does not depend on
a numerical extrapolation in $m$.

Let $F$ be one such polynomial of bidegree $(a,g)$, $a+g=5$.
For a row polynomial, express it uniquely as
\begin{equation}\label{R20:rowbasis}
 F=\sum_{j=1}^{a+1} Q^{\rm row}_{a,j}(x,y) C_j(z,w),\qquad
 Q^{\rm row}_{a,j}=
 \prod_{u=j}^{a}(x-Kq^uy)\prod_{u=1}^{j-1}(x-q^uy).
\end{equation}
Define the binary quintic
\begin{equation}\label{R20:rowtransform}
 \mathcal R(F)=\sum_{j=1}^{a+1}X^j C_j(z,q^{-j}w)
       \prod_{u=5-j}^{3}(z-q^uw)
       \prod_{u=-1}^{a-1-j}(z-Kq^uw).
\end{equation}
As everywhere here, a product with lower endpoint greater than its
upper endpoint is one. For a column polynomial use
\begin{equation}\label{R20:colbasis}
 \begin{split}
 F&=\sum_{j=0}^{a}Q^{\rm col}_{a,j}(x,y)C_j(z,w),\\
 Q^{\rm col}_{a,j}&=
 \prod_{u=5-j}^{4}(x-q^uy)
 \prod_{u=0}^{a-1-j}(x-K^{-1}q^{5-a+u}y),
 \end{split}
\end{equation}
and define
\begin{equation}\label{R20:coltransform}
 \mathcal S(F)=\sum_{j=0}^{a}X^{-j}C_j(z,Kq^{j-5}w)
       \prod_{u=-a+j}^{-1}(z-q^uw)
       \prod_{u=0}^{j-1}(z-Kq^{-5+u}w).
\end{equation}
The coefficient matrices of the $Q$ bases are generically invertible.
Their nonzero determinants are also included in the certificate below,
so no problematic specialization is silently inverted.

The row matrix consists of coefficient columns, in the order
$z^5,z^4w,\ldots,w^5$, of
\[
 z\prod_{u=0}^3v_u,\quad w\prod_{u=0}^3v_u,
 \quad\mathcal R(R_0),\ldots,\mathcal R(R_3).
\]
The column matrix consists of the corresponding columns of
\[
 z\prod_{u=-5}^{-2}(z-Kq^uw),\quad
 w\prod_{u=-5}^{-2}(z-Kq^uw),
 \quad\mathcal S(S_0),\ldots,\mathcal S(S_3).
\]
At $(q,K,X)=(6,-1,-1)$ over $\mathbb F_{41}$, these matrices are
\begin{equation}\label{R20:matrices}
 \begin{split}
 M_R&=\begin{pmatrix}
 1&0&36&1&21&34\\28&1&24&28&24&4\\
 34&28&3&28&35&22\\21&34&34&6&27&27\\
 39&21&0&33&24&34\\0&39&26&27&22&6
 \end{pmatrix},\\
 M_C&=\begin{pmatrix}
 1&0&30&28&31&1\\2&1&33&10&24&8\\
 19&2&1&9&13&0\\34&19&35&4&27&5\\
 2&34&13&36&15&4\\0&2&32&3&11&27
 \end{pmatrix}.
 \end{split}
\end{equation}
Their determinants are respectively $16$ and $15$ modulo $41$.
The four row interpolation determinants are $29,1,29,38$;
the four column interpolation determinants are $25,6,32,32$.
All entries are obtained by the products and linear systems
\eqref{R20:rowbasis}--\eqref{R20:coltransform}, which specify a
small reproducible exact certificate.

We record explicit degree bounds in order to turn this certificate
into a uniform statement. Let $d_R$ be the determinant of a row
coefficient matrix in \eqref{R20:rowbasis}, and let $d_C$ be the
determinant of the coefficient matrix of
$K^a Q^{\rm col}_{a,j}$ in \eqref{R20:colbasis}.
For all eight forms, $0\leq a\leq3$. Cramer's rule shows that
the coefficients of
\[
 q^{60}X^5d_R\mathcal R(F),\qquad
 q^{60}X^5d_C\mathcal S(F)
\]
are polynomials in $q,K,X$. Each has degree at most
$150,20,10$ in those three variables separately. Here are direct
bounds verifying the clearing and the stated estimates.
Every coefficient of a local form has $q$-degree at most $20$.
A row interpolation entry has $q$-degree at most $a(a+1)/2$
and $K$-degree at most $a$; its adjugate entry has degrees at most
$18,9$. The substitutions and additional products in
\eqref{R20:rowtransform} have negative $q$ exponent no smaller
than $-10$ and positive contribution at most $7$; after clearing,
the $q$-degree is at most $105$, the $K$-degree at most $12$,
and the $X$-degree at most $9$.
For the column basis, multiplying by $K^a$ gives entries of degrees
at most $4a,a$. Its adjugate has degrees at most $36,9$.
The solution acquires an additional factor $K^a$, and
\eqref{R20:coltransform} has negative $q$ exponent no smaller than
$-31$. After clearing its degrees are at most $116,17,5$.
These are bounded by $150,20,10$ as asserted.

Multiply the first two columns of each matrix by $q^{25}$,
and clear each of its four other columns as just described.
The two base columns now have degrees at most $31,4,0$.
Let $\Delta_R,\Delta_C$ be the two resulting polynomial determinants,
and define
\[
 \mathscr F(q,K,X)=\Delta_R\Delta_C
                  \prod_{\text{four row forms}}d_R
                  \prod_{\text{four column forms}}d_C.
\]
Each interpolation determinant has degrees at most $48,12,0$.
Consequently
\begin{equation}\label{R20:degrees}
 \deg_q\mathscr F\leq2000,\qquad
 \deg_K\mathscr F\leq300,\qquad
 \deg_X\mathscr F\leq80.
\end{equation}
For example the sharper bounds from the preceding estimates are
$1708,272,80$. At the certificate specialization all eight
interpolation determinants have product $20$ modulo $41$;
$K^{a(a+1)}=1$ there, so the column clearing does not change their
listed values. Thus
\[
 \mathscr F(6,-1,-1)
 =6^{580}(-1)^{40}\,20^2\,16\,15=22\pmod{41}.
\]
In particular $\mathscr F$ is a nonzero integer polynomial.

\paragraph{3. Why these matrices are the required Fourier quotients.}
Use the cyclic action $\tau:(y,w)\mapsto(qy,qw)$ of R-16 and
write $A(x,y)=\prod_{u=k}^{m-1}\ell_u$.
All its coefficients $A_r$, $0\leq r\leq m-k$, are nonzero by
the geometric-product formula already proved there. In character
$t$, divide the coefficient with $y$-degree $r$ and $w$-degree
$s$ by $A_r$; here $r+s=t$.

For a row form of $(x,y)$ degree $a$, the outside factor is
$\ell_0\prod_{u=k+a+1}^{m-1}\ell_u$ in those variables.
Multiplication by each $Q^{\rm row}_{a,j}$ in
\eqref{R20:rowbasis} completes exactly $A(x,q^jy)$.
Dividing its coefficient by $A_r$ gives $q^{jr}=X^jq^{-js}$.
The remaining $(z,w)$ factors are
$C_j(z,q^{-j}w)\prod_{u=5-j}^{k+a-1-j}v_u$.
Their common product is
\[
 H_R=\prod_{u=4}^{k-2}v_u,
\]
and the residual quintic is precisely
\eqref{R20:rowtransform}.

For a column form, the outside $(x,y)$ factor is
$\prod_{u=5}^{c-1}\ell_u$, with $c=m+5-a-k$.
The $j$th basis term in \eqref{R20:colbasis} completes
$A(x,q^{5-k-j}y)$. After coefficient division and multiplication
by the nonzero character scalar $X^{k-5}$, its factor is
$X^{-j}q^{(k-5+j)s}$.
Shifting the outside $(z,w)$ roots $c,\ldots,m-1$ by
$k-5+j$ gives indices $-a+j,\ldots,k-6+j$ modulo $m$.
Their common product is
\[
 H_C=\prod_{u=0}^{k-6}v_u,
\]
and the residual quintic is exactly
\eqref{R20:coltransform}. This checks both shifts and both
normalizations, including their dependence on $k,t$.

Suppose $9\leq k,t\leq m-9$ and
$\mathscr F(q,q^k,q^t)\ne0$. The two matrices then show that
base plus row columns span $H_R\operatorname{Sym}^5\mathbb C^2$,
and base plus column columns span
$H_C\operatorname{Sym}^5\mathbb C^2$ before restriction to the
allowed coefficient interval. Their common factor is
$H=\prod_{u=4}^{k-6}v_u$. After dividing by $H$, the first space is
\[
 \left(\prod_{u=k-5}^{k-2}v_u\right)
                 \operatorname{Sym}^5\mathbb C^2,
\]
and the second is
\[
 \left(\prod_{u=0}^{3}v_u\right)
                 \operatorname{Sym}^5\mathbb C^2.
\]
The two displayed degree-four factors are coprime: their indices
are distinct and lie below $m$. Their intersection consists exactly
of their product times $\operatorname{Sym}^1\mathbb C^2$,
which is the two-dimensional base space. The joint span is therefore
all $H\operatorname{Sym}^9\mathbb C^2$, of dimension ten.
This proves joint independence; separate row and column ranks alone
would not suffice.

There are ten consecutive allowed coefficient indices
$b,\ldots,b+9$ with $0\leq b\leq k-9$.
Multiplication by $H$ and restriction to them has determinant
\[
 \det(H_{b+u-v})_{0\leq u,v\leq9}
 =\pm s_{(10^b)}(q^4,q^5,\ldots,q^{k-6})\ne0.
\]
This is the same dual Jacobi--Trudi argument as R-16/R-19
(the cited unconditional identity R-B15). There are $N=k-9$
variables; the numerator exponents are distinct in $[0,N+9]=[0,k]$,
strictly below $m$, so the Weyl alternant is a nonzero Vandermonde.
For $N=0$ the coefficient matrix is the identity.
Consequently the eight normal classes remain independent modulo
the entire base character space after the required coefficient
restriction.

\paragraph{4. Only finitely many modes are lost.}
Choose a nonzero coefficient polynomial in $q$ of the nonzero
polynomial $\mathscr F(q,K,X)$. Its degree is at most $2000$.
For every positive integer $m$,
\[
 \varphi(m)^2/m
 =\prod_{p^e\Vert m}p^{e-2}(p-1)^2\geq\tfrac12:
\]
each odd-prime factor is at least one, and the possible factor
from $p=2$ is at least $1/2$.
For $m\geq M_0$ there are therefore more than $2000$ primitive
$m$th roots. Choose one at which that coefficient polynomial does
not vanish. This uses ordinary root counting and gives
$\mathscr F(q,K,X)\not\equiv0$; it does not incorrectly replace
$\varphi(m)$ by $m$.

For this $q$, at most $300$ of the distinct values $K=q^k$ can
make $\mathscr F(q,K,X)$ the zero polynomial in $X$: choose one
nonzero coefficient polynomial in $K$ and use
\eqref{R20:degrees}. Among $k=9,\ldots,m-9$ there are consequently
at least $m-317$ good bidegrees. For each good $k$, at most $80$
of the distinct values $X=q^t$ are zeros. Thus at least $m-97$
characters in $t=9,\ldots,m-9$ have all eight independent classes.

Consider translates of the eight leading vectors at the $h_m$
block starts $a=5j$, $0\leq j<h_m$. A relation in a good bidegree,
projected onto a good character, gives eight separate equations
\[
 \sum_{j=0}^{h_m-1}c_j(q^{5t})^j=0.
\]
The good characters contain at least $h_m$ distinct values of
$q^{5t}$: every value has at most $\gcd(5,m)\leq5$ preimages,
and $h_m=\lfloor(m-97)/5\rfloor$.
Each equation system contains an ordinary invertible Vandermonde
matrix. Hence the $8h_m$ leading normal classes in each good
bidegree are independent modulo $W_0$.

The original polynomial limits need not be homogeneous for this
bigrading. Nevertheless they are independent modulo $W_0$ as well.
In any relation, take the largest $k$ among vectors with nonzero
coefficients. Their components of bidegree $(m-k,k)$ are precisely
the independent leading classes just proved; every other component
has smaller $(z,w)$ degree. Because $W_0$ is bigraded, its contribution
can be removed in that quotient, forcing all coefficients at this
$k$ to vanish. Descending in $k$ proves independence of every
selected vector. This justifies the leading-component step without
assuming that a projection of the deformed differential image lies
in that image.

Finally, the eight sets of polynomial families all belong, for
nonzero $\varepsilon$, to the differential image of the same
simultaneous pencil. Indeed \eqref{R20:relation} times its outside
permanent is a linear-form combination of global complementary
permanents; division by $\varepsilon^2$ is polynomial by
\eqref{R20:secondorder}. Their limits and $2m^2+2$ base columns
have a nonzero augmented minor at zero. This minor stays nonzero
for a common generic nonzero $\varepsilon$. Thus the dimension
bound \eqref{R20:dimension} follows from the actual differential
rank, with no addition of ranks from different pencils.

\paragraph{5. The uniform padding consequence.}
The general criterion \eqref{R20:allrange} is R-14, with G-10
giving its effective coordinate-degree threshold. Moreover
\[
 h_m\geq(m-101)/5,
\]
so
\begin{align*}
 2m^2+J_m+2-\frac{32m^2}{9}
 &\geq \frac{2m^2}{45}-\frac{3344m}{5}
                  +\frac{256136}{5}+2\;>0
 \qquad(m\geq M_0).
\end{align*}
The last inequality already holds for $m>15048$ after dropping
the positive constant. Thus every integer $m<n<4m/3$ satisfies
\eqref{R20:allrange}. Its strict padding inequality is exactly
the stable range required by G-11, yielding the stated explicit
rectangles. G-12 excludes rectangular gaps for every $n\geq4m/3$.
Finally $J_m=(8/5)m^2+O(m)$, so the full dimension criterion reaches
every fixed ratio below $\sqrt{9/5}=3/\sqrt5$; above $4/3$, the
resulting four-row gaps must be nonrectangular.
\end{proof}


\subsection{R-21: arithmetic rank recovery and a fixed-dimension algorithm}

\paragraph{Statement R-21.}
The integer Gram model admits an explicit spectral separation from
zero. Moreover, for each fixed ambient dimension $N$, every orbit-closure
multiplicity $m_\lambda(f)$ of a rational degree-$n$ form can be computed
exactly in time polynomial in the numerical degree $D=nd$ and the
coefficient bit length. This is an elementary consequence of a compressed
coefficient kernel and weight extraction. The exponent of the polynomial
depends on $N$; this assertion is not polynomial in $\log d$ and does not
give a uniform polynomial-time algorithm when $N$ grows.

\paragraph{1. An integral spectral bound.}
Let $A$ be an $a\times b$ integer matrix of positive rank $r$, with
entries of absolute value at most $H$, and put $G=AA^{\mathsf T}$,
$T=\operatorname{tr}G\leq abH^2$. Then
\begin{equation}\label{R21:gap}
 \lambda_{\min}^{+}(G)\geq
 \begin{cases}
 1,&r=1,\\[1mm]
 \bigl((r-1)/T\bigr)^{r-1},&r>1.
 \end{cases}
\end{equation}
More generally the same bound holds for any nonzero positive semidefinite
integer symmetric matrix, with its own rank and trace.

\begin{proof}
If the positive eigenvalues are $\eta_1\leq\cdots\leq\eta_r$, their
product is the elementary symmetric coefficient $e_r(G)$, a positive
integer. For a Gram matrix it is also a sum of squared integer minors
by Cauchy--Binet. Thus $\prod_i\eta_i\geq1$.
If $r=1$ this proves the claim. Otherwise the arithmetic--geometric
mean inequality gives
\[
 \prod_{i=2}^r\eta_i\leq
 \left(\frac{\sum_{i=2}^r\eta_i}{r-1}\right)^{r-1}
 \leq\left(\frac{T}{r-1}\right)^{r-1},
\]
which proves \eqref{R21:gap}. The proof for an integer positive
semidefinite matrix uses the same integer characteristic coefficient.
\end{proof}

The exponential dependence on matrix dimension can occur even at fixed
entry height. If $S$ is the nilpotent superdiagonal shift of size $r$
and $A=I-HS$ with integer $H\geq2$, then
$A^{-1}=\sum_{j=0}^{r-1}H^jS^j$. Its top-right entry is $H^{r-1}$,
so $\|A^{-1}\|_2\geq H^{r-1}$ and
$\lambda_{\min}(AA^{\mathsf T})\leq H^{-2(r-1)}$.
This is a concrete limitation on dimension-independent spectral
separation from integrality alone.

For example, let $A_\lambda$ be the integral matrix in R-06 constructed
from the tensor $n!f$, where the original polynomial has integer monomial
coefficients. Set $c_\lambda=f^\lambda/D!$. Its projected Gaussian
covariance for the original, unscaled symmetric tensor is
\begin{equation}\label{R21:gaussiantransport}
 \mathcal G_\lambda(f)=
 \frac{c_\lambda^2}{(n!)^{2d}}
 A_\lambda\operatorname{diag}(\alpha!)A_\lambda^{\mathsf T}.
\end{equation}
When $A_\lambda\ne0$, every $\alpha!\geq1$, so its least positive eigenvalue is at least
$c_\lambda^2/(n!)^{2d}$ times the bound in \eqref{R21:gap} for
$A_\lambda A_\lambda^{\mathsf T}$. Both matrices have the same image,
so the comparison is made on their common support. The factorials in
this statement are necessary; the passage from an integral polynomial
to an integral symmetric tensor cannot be omitted.
The explicit matrix dimensions and height in R-06 show that this
bound alone does not yield polynomial precision in growing GCT size.

\paragraph{2. A compressed integral coefficient kernel.}
Assume $f=\sum_{|\beta|=n}f_\beta x^\beta$ has integer coefficients;
rational inputs can be cleared of a common denominator without changing
the orbit span. Let $f_{\rm ten}\in(\mathbb C^N)^{\otimes n}$ be the
usual symmetric-tensor embedding and set
$v=(n!f_{\rm ten})^{\otimes d}$.
For matrices of variables $Z,W$, write
\[
 Z\cdot f=\sum_{|\beta|=n}g_\beta(Z)x^\beta,
 \qquad
 \Phi(Z)=Z^{\otimes D}v=\sum_{|\alpha|=D}Z^\alpha c_\alpha.
\]
There are
\[
 M=\binom{N^2+D-1}{D}
\]
column labels $\alpha$. Form the polynomial
\begin{equation}\label{R21:kernel}
 \mathcal K_f(Z,W)=
 \left(n!\sum_{|\beta|=n}\beta!\,
                     g_\beta(Z)g_\beta(W)\right)^d
       =\sum_{\alpha,\gamma}B_{\alpha\gamma}Z^\alpha W^\gamma.
\end{equation}
Then $B=C_f^{\mathsf T}C_f$ is an $M\times M$ positive semidefinite
integer matrix. Thus its kernel coefficients replace the exponentially
tall tensor-coordinate matrix without constructing its $N^D$ rows.

\begin{proof}
For a composition $\beta$ of $n$, the symmetric-tensor embedding of
the monomial $x^\beta$ is
$(\beta!/n!)$ times the sum of all distinct tensor words with
composition $\beta$. There are $n!/\beta!$ such words.
Consequently the coefficient of each such word in
$Z^{\otimes n}(n!f_{\rm ten})$ is $\beta!g_\beta(Z)$, and its
bilinear tensor inner product with the analogous expression in $W$ is
\[
 \sum_{|\beta|=n}\frac{n!}{\beta!}
                 (\beta!g_\beta(Z))(\beta!g_\beta(W))
 =n!\sum_{|\beta|=n}\beta!g_\beta(Z)g_\beta(W).
\]
Taking $d$ tensor powers gives \eqref{R21:kernel}. Comparing coefficients
gives $B_{\alpha\gamma}=c_\alpha^{\mathsf T}c_\gamma$.
The vectors are real and integral; the resulting real Gram matrix is
also positive semidefinite for the corresponding Hermitian inner product.
\end{proof}

For $\alpha=(\alpha_{ij})$, define its row-sum weight by
\[
 \operatorname{wt}(\alpha)_i=\sum_{j=1}^N\alpha_{ij}.
\]
The matrix $B$ is block diagonal by these weights. Denote the block
of weight $\mu$ by $B_\mu$ and set $r_\mu=\operatorname{rank}B_\mu$,
with $r_\mu=0$ if a coordinate of $\mu$ is negative or the block is
absent. For $\delta=(N-1,N-2,\ldots,0)$, the exact formula is
\begin{equation}\label{R21:weightformula}
 m_\lambda(f)=
 \sum_{\sigma\in S_N}\operatorname{sgn}(\sigma)
              r_{\lambda+\delta-\sigma\delta},
 \qquad \lambda\vdash D,\quad\ell(\lambda)\leq N.
\end{equation}
This is a signed Weyl extraction from positive Gram ranks, not a new
cancellation-free tableau rule.

\begin{proof}
For a diagonal matrix $t$, the identity
$\Phi(tZ)=t^{\otimes D}\Phi(Z)$ shows that $c_\alpha$ has weight
$\operatorname{wt}(\alpha)$. Distinct tensor-weight spaces are
orthogonal in the word basis, proving the block decomposition.
As in R-06, the coefficient span is exactly the $GL_N$ orbit span:
a functional vanishes on every coefficient if and only if its
polynomial evaluation vanishes on all matrices, equivalently on the
Zariski-dense invertible matrices. Hence $r_\mu$ is exactly the
dimension of the $\mu$ weight space of this orbit span.

Its character is $\sum_\mu r_\mu z^\mu$.
Multiplication by the Weyl alternant
$\sum_{\sigma\in S_N}\operatorname{sgn}(\sigma)z^{\sigma\delta}$
turns $s_\nu(z)$ into the alternant with exponent $\nu+\delta$.
The coefficient of $z^{\lambda+\delta}$ is therefore $m_\lambda(f)$:
the strictly decreasing dominant exponent occurs in precisely that
alternant, with coefficient one. Expanding this coefficient proves
\eqref{R21:weightformula}. The orbit-span multiplicities are the
coordinate-ring multiplicities under the consistent dual convention
already fixed in R-05.
\end{proof}

\paragraph{3. Exact cost and sufficient approximation precision.}
Let $H=\max\{1,|f_\beta|\}$ and put
\[
 H_B=N^{3D}(n!H)^{2d},\qquad T_0=MH_B.
\]
Every entry of $B$ has absolute value at most $H_B$.
Indeed each tensor coefficient has absolute value at most
$N^D(n!H)^d$, and each Gram entry sums $N^D$ products of two such
coefficients. Thus $\operatorname{tr}B_\mu\leq T_0$ for every block.
The spectral lemma, without knowing the block rank, gives
\begin{equation}\label{R21:uniformgap}
 \lambda_{\min}^{+}(B_\mu)\geq
 \eta_0:=T_0^{-(M-1)}
\end{equation}
for every nonzero block. To see this, use rank at most $M$,
$T_0\geq1$, and discard the factor $(r-1)^{r-1}\geq1$ from
\eqref{R21:gap}; a rank-one positive eigenvalue is at least one.

If a Hermitian approximation to each block has entrywise error
less than $\eta_0/(4M)$, its operator error is less than
$\eta_0/4$. The min--max principle then places each eigenvalue
within that distance of its exact value. Counting the approximate
eigenvalues above $\eta_0/2$ recovers the exact rank. For example,
entry precision of
\begin{equation}\label{R21:bits}
 p=\left\lceil (M-1)\log_2T_0+\log_2(8M)\right\rceil
\end{equation}
absolute entry error at most $2^{-p}$ suffices, with room for an
additional eigenvalue-approximation error at most $\eta_0/8$. For this integral coefficient matrix there is also a simpler exact
recovery route: absolute entry error below $1/2$ permits rounding each
entry to its exact integer, followed by exact elimination. The smaller
error above is needed only for the stated direct spectral threshold;
it must not be confused with a necessary precision for integer recovery.
The spectral statement is an explicit conditioning bound; it does not by itself bound the cost of a
random-matrix integration oracle.

For fixed $N$, $M$ is polynomial in $D$, and
\[
 \log H_B=O\bigl(D\log N+D\log n+d\log H\bigr).
\]
The $g_\beta$ and the kernel in \eqref{R21:kernel} can be expanded
by integer polynomial arithmetic in a fixed number of variables.
Their dense coefficient arrays, at degrees at most $D$ in each
matrix alphabet, have polynomial size in $D$.
Their coefficient bit lengths are bounded by the same elementary
factorial and multinomial estimates. Thus all kernel entries can
be computed in polynomial bit complexity for fixed $N$.
Exact rational elimination computes each required block rank in
polynomial bit complexity as well: intermediate reduced fractions
are ratios of minors, whose sizes are bounded by the factorial
expansion bound $M!H_B^M$. Formula \eqref{R21:weightformula} needs
at most $N!$ block ranks for a given label. This proves the claimed
fixed-dimension algorithm; \eqref{R21:bits} is polynomial in the
same numerical parameters.

There are two relevant limits to this conclusion. In GCT the ambient
dimension is $N=n^2$, so neither the size exponent nor the precision
bound is uniform polynomial in matrix size. Also, an orbit of one
fixed quaternary form is not the entire image of all four-variable
permanent pencils. Applying this algorithm to one specialization
does not compute the multiplicities of $\mathcal P_{m,4}$.

\paragraph{The slope-method gate.}
The supplied Alp\"oge note, step~13 on printed p.~4 (G-B10), names
the Bost slope inequality as a step. That step itself gives no integral-metric construction or
degree-dependent constants for the coefficient matrices above. A transfer to this rank problem would require such
a construction and bounds strong enough to control the displayed
dimensions and heights. The present spectral estimate supplies a
self-contained arithmetic bound and an exact restricted algorithm;
it neither uses nor challenges the stipulated RH/Dirichlet-GRH input.
\subsection{R-22: arbitrary staircase size reaches the four-variable dimension ceiling}

\paragraph{Statement R-22.}
Let $b\geq3$ and $m\geq8b$ be integers, and define
\[
 T=m-4b+3,\qquad h=\left\lfloor\frac{m-6b+3}{b}\right\rfloor,
 \qquad J_{m,b}=(2b-2)hT.
\]
Then the quaternary permanent projection variety satisfies
\begin{equation}\label{R22:rank}
 c_m:=\dim\mathcal P_{m,4}\ \geq\ 2m^2+2+J_{m,b}.
\end{equation}
The proof uses one simultaneous deformation of the cyclic pencil.
Its untwisted motif parameters lie in the integer grid
$\{0,\ldots,4(b-1)T\}$; the translated copies have the prescribed
root-of-unity twists. A grid choice is allowed to depend on $m$ and on
the chosen primitive $m$th root of unity. No ranks from
different pencils are added.

In particular, for every $m\geq99$,
\begin{equation}\label{R22:asymptotic}
 4m^2-20m^{3/2}+2\ \leq\ c_m\ \leq\ 4m^2-2m+2.
\end{equation}
Thus $c_m=4m^2-O(m^{3/2})$. The upper bound is the row--column torus
bound proved in R-13, and an equality formula for $c_m$ is not asserted.

For any integers $n>m$ such that
\begin{equation}\label{R22:padding}
 2n^2\leq2m^2+J_{m,b}+2,
\end{equation}
the padded image has dimension greater than the determinant image.
Consequently R-14/G-10 give exactly-four-row multiplicity gaps in every
coordinate degree beyond the explicit G-10 threshold. The specified
stretched labels and effective bounds of G-17 also apply.
For every fixed real $1<\rho<\sqrt2$, this includes every integer
$m<n\leq\rho m$ as soon as
\begin{equation}\label{R22:ratio_threshold}
 m\geq\max\left\{99,
       \left\lceil\frac{100}{(2-\rho^2)^2}\right\rceil\right\}.
\end{equation}
In particular the eventual rectangular boundary of R-20 holds already
for every $m\geq2025$: rectangular gaps exist for every integer
$m<n<4m/3$, and G-12 excludes them for $n\geq4m/3$.
The additional range above $4m/3$ consists of nonrectangular gaps.

\begin{proof}
We give a uniform local calculation, prove its two Fourier rank
conditions, and then arrange that one integer choice works for all
interior bidegrees of the given matrix size.

\paragraph{1. The typed staircase and its simultaneous supports.}
Fix any primitive $m$th root $q$, and write
\[
 \ell_i=x-q^iy,\qquad v_i=z-q^iw.
\]
The cyclic pencil has entries $\ell_i$ at $(i,i)$ and $v_i$ at
$(i,i+1)$, with indices modulo $m$.
A block of $b$ rows, initially rows $0,\ldots,b-1$, has columns
$0,\ldots,b$. In addition to its diagonal and superdiagonal, allow
every entry $(i,j)$ except entries in column zero and the entry
$(0,b)$. Thus the permitted additional positions are
\[
 0\leq i<b,\quad1\leq j\leq b,\quad
 j\ne i,i+1,\quad(i,j)\ne(0,b).
\]
There are $b(b-2)$ such positions. Set their entries to
$\varepsilon C_{ij}$, where
\begin{equation}\label{R22:types}
 C_{ij}=\begin{cases}
 u_{ij}x+w_{ij}y,&j>i+1,\\
 u_{ij}z+w_{ij}w,&j<i.
 \end{cases}
\end{equation}
The symbol $w_{ij}$ denotes a scalar coefficient, distinct from the
coordinate $w$. All $2b(b-2)$ coefficients are initially independent
parameters, collectively denoted $\theta$.

Place copies of this same block at starts $a=bj$, $0\leq j<h$,
replacing $(y,w)$ in its additional forms by $(q^ay,q^aw)$.
All block rows are disjoint, since $bh\leq m-6b+3<m$.
For each block, its bottom $b-1$ rows have support only in its
$b$ columns numbered $1,\ldots,b$ relative to the start.
Its middle $b-1$ columns have support only in its $b$ rows numbered
$0,\ldots,b-1$. A neighboring block shares only a boundary column,
so both support assertions remain true with all blocks present.

These supports give exact cofactor factorizations. If an external row
is deleted, and one of the $b$ supported columns is deleted, all
bottom $b-1$ supported rows must occupy the remaining $b-1$ supported
columns in every matching. The complementary permanent is the product
of their local permanent minor and a common outside permanent.
The transposed assertion holds when an external column and one of
the $b$ supported rows are deleted. Other blocks can affect the outside
factor but cannot alter either factorization.

\paragraph{2. The second-order local defect for arbitrary size.}
Put $r=b-1$. Consider an $r\times(r+1)$ matrix $B(\varepsilon)$
with diagonal entries $a_i$, superdiagonal entries $b_i$, and
other entries $\varepsilon c_{ij}$, with $0\leq i<r$.
Form $\widetilde B$ by signing its additional entries with
\[
 \sigma_{ij}=\begin{cases}
 (-1)^{j-i-1},&j>i+1,\\
 (-1)^{i-j},&j<i.
 \end{cases}
\]
Let $P_j$ be the permanent minor deleting column $j$, and $D_j$
the corresponding determinant minor of $\widetilde B$.
There is one order-zero matching. A matching with one additional
edge has its other edges forced along the bidiagonal path; its
permutation sign is precisely the indicated $\sigma_{ij}$.
Hence
\begin{equation}\label{R22:local_defect}
 P_j-D_j=\varepsilon^2E_j+O(\varepsilon^3).
\end{equation}

We need the endpoint identities $E_0=E_r=0$ for all $r$.
For the minor deleting column $r$, a matching with no additional
forward edge is an upper-Hessenberg permutation. Each nontrivial
cycle is a consecutive interval, whose superdiagonal edges end in
one backward edge. Its permutation sign equals the sign assigned
to that backward edge. Two forward additional edges and no backward
edge cannot form a permutation. If there is exactly one forward
edge $i\mapsto j$, $j>i+1$, and one backward edge $h\mapsto l$,
there is a single nontrivial cycle, necessarily with
$l\leq i<j\leq h$. It follows superdiagonal edges from $l$ to $i$,
jumps to $j$, follows them to $h$, and returns to $l$.
Its length is $(i-l)+(h-j)+2$. The difference between its sign
exponent and the sum of the two assigned exponents is even, since
\[
 (j-i-1)+(h-l)-\big((i-l)+(h-j)+1\big)=2(j-i-1).
\]
Thus the endpoint minor agrees through order two.
Deleting column zero and transposing gives the same square
bidiagonal situation, with diagonal $b_i$ and superdiagonal
$a_{i+1}$. Under this operation a forward edge becomes a backward
edge with the same assigned sign, and conversely. This proves
\begin{equation}\label{R22:endpoints}
 E_0=E_r=0.
\end{equation}

Assume now that $a_i$ and the additional forward entries are
linear in $(x,y)$, while $b_i$ and the additional backward entries
are linear in $(z,w)$. The highest $(z,w)$ degree of $E_j$ is
$r-1$. Indeed, if there is no additional forward edge, the same
upper-Hessenberg argument gives equality of permanent and signed
determinant for every minor: the rows before the deleted column
form the upper-Hessenberg block, and all later superdiagonal
entries are forced. Therefore terms with no $(x,y)$ factor, and
terms with a diagonal $(x,y)$ factor and two backward edges, cancel.
At $(x,y)$ degree one the only remaining possibility uses one
forward edge, one backward edge, and otherwise superdiagonal edges.

For the minor deleting column $j$, compare its target columns with
the all-superdiagonal target set $\{1,\ldots,r\}$.
The target zero must come from the backward edge. Comparing the
two removed targets with the two added targets then forces the
additional edges to be
\[
 (j-1,h+1),\qquad(h,0),\qquad j\leq h\leq r-1.
\]
The minor permutation has the cycle
$(0,1,\ldots,j-1,h)$, of sign $(-1)^j$; the product of assigned
edge signs has opposite sign. Consequently
\begin{equation}\label{R22:Etop}
 [E_j]_{\deg_{z,w}=r-1}
 =2\sum_{h=j}^{r-1}c_{j-1,h+1}c_{h,0}
           \prod_{\substack{0\leq t<r\\t\ne j-1,h}}b_t
 \quad(1\leq j\leq r-1).
\end{equation}
This proves the formula for every block size, without extrapolating
a finite list of local minors.

The determinant row relation yields
\begin{equation}\label{R22:relation}
 \sum_{j=0}^{r}(-1)^{i+j}\widetilde B_{ij}P_j
 =\varepsilon^2(a_iE_i-b_iE_{i+1})+O(\varepsilon^3).
\end{equation}
For $i<r-1$, the highest component, divided by two, is
\[
 -\sum_{h=i+1}^{r-1}c_{i,h+1}c_{h,0}
                          \prod_{t\ne h}b_t,
\]
of bidegree $(1,r)$. For $i=r-1$ it is
\[
 a_{r-1}c_{r-2,r}c_{r-1,0}\prod_{t=0}^{r-3}b_t,
\]
of bidegree $(2,r-1)$.

In the original $b$-row block, denote these row forms by
$R_0,\ldots,R_{r-1}$. Explicitly,
\begin{align}
 R_i&=-\sum_{h=i+1}^{r-1}C_{i+1,h+2}C_{h+1,1}
                 \prod_{\substack{0\leq t<r\\t\ne h}}v_{t+1}
                     &&(0\leq i<r-1),\label{R22:Rforms}\\
 R_{r-1}&=\ell_r C_{r-1,r+1}C_{r,1}
                                      \prod_{t=1}^{r-2}v_t.\notag
\end{align}
The transposed supported-column calculation gives forms
$S_0,\ldots,S_{r-1}$:
\begin{align}
 S_0&=-\ell_1 C_{0,2}C_{r,1}\prod_{u=2}^{r-1}v_u,
                                                     \label{R22:Sforms}\\
 S_i&=\sum_{p=0}^{i-1}C_{p,i+1}C_{r,p+1}
                    \prod_{\substack{0\leq u<r\\u\ne p}}v_u
                    &&(1\leq i<r).\notag
\end{align}
Their $(x,y)$ degrees are respectively
\begin{equation}\label{R22:grades}
 (1,\ldots,1,2)\quad\hbox{and}\quad(2,1,\ldots,1).
\end{equation}
All their coefficients are homogeneous quadratic polynomials in
$\theta$. These displayed forms are highest components; the actual
limits from \eqref{R22:relation} can also have smaller $(z,w)$ degree.
We retain those lower components until the final independence argument.

\paragraph{3. The Fourier matrices and their exact meanings.}
Fix
\[
 2b-1\leq k,t\leq m-2b+1,\qquad K=q^k,\quad X=q^t.
\]
At $\varepsilon=0$, the common outside factors for the block at zero
are exactly
\begin{equation}\label{R22:outside}
 T_s^R=\ell_0\prod_{i=b}^{s-1}v_i\prod_{i=s+1}^{m-1}\ell_i,
 \qquad
 T_c^C=\prod_{i=b}^{c-1}\ell_i\prod_{i=c}^{m-1}v_i.
\end{equation}
The unique cyclic matching proves these formulas. For a local form
of bidegree $(a,g)$, $a+g=b$, take
\begin{equation}\label{R22:alignment}
 s=k+b-g=k+a,\qquad c=m+g-k.
\end{equation}
These indices are external to the corresponding supported rows or
columns throughout the stated range. The resulting leading products
have bidegree $(m-k,k)$.

Set $A(x,y)=\prod_{u=k}^{m-1}\ell_u$ and let $A_i$ be its coefficient
of $x^{m-k-i}y^i$. All $A_i$ are nonzero: the finite geometric-product
formula for these coefficients has numerator and denominator factors
$1-q^j$ with $0<j<m$ (the same formula used in R-16).
In cyclic character $t$, divide the coefficient with $y$-degree $i$
and $w$-degree $j$, $i+j=t$, by $A_i$.
Under this normalization the entire cyclic base image $W_0$ has the
two-dimensional binary coefficient space
\begin{equation}\label{R22:base}
 \operatorname{span}\left\{
 z\prod_{u=0}^{k-2}v_u,
 w\prod_{u=0}^{k-2}v_u\right\}.
\end{equation}
This is the base-character description proved in R-16; the full base
image has dimension $2m^2+2$ by R-15/R-17.

For a row form $F$ of $(x,y)$ degree $a\in\{1,2\}$, write
\begin{align}
 F&=\sum_{j=1}^{a+1}Q^R_{a,j}(x,y)C_j(z,w),\notag\\
 Q^R_{a,j}&=\prod_{u=j}^{a}(x-Kq^uy)
                       \prod_{u=1}^{j-1}(x-q^uy),\label{R22:rowbasis}\\
 \mathcal R_b(F)&=\sum_{j=1}^{a+1}X^j C_j(z,q^{-j}w)
       \prod_{u=b-j}^{b-2}v_u
       \prod_{u=-1}^{a-1-j}(z-Kq^uw).\label{R22:rowtransform}
\end{align}
For a column form use
\begin{align}
 F&=\sum_{j=0}^{a}Q^C_{a,j}(x,y)C_j(z,w),\notag\\
 Q^C_{a,j}&=\prod_{u=b-j}^{b-1}(x-q^uy)
               \prod_{u=0}^{a-1-j}(x-K^{-1}q^{b-a+u}y),
                                                     \label{R22:colbasis}\\
 \mathcal S_b(F)&=\sum_{j=0}^{a}X^{-j}C_j(z,Kq^{j-b}w)
       \prod_{u=-a+j}^{-1}v_u
       \prod_{u=0}^{j-1}(z-Kq^{-b+u}w).\label{R22:coltransform}
\end{align}
An empty product is one. All displayed bases are invertible for the
specified $q,k$. For the row bases, evaluation at $x=qy$, and then
at $x=q^2y$ when $a=2$, gives triangular independence; the possible
nonconstant vanishing factors are $q-1$, $K-1$, and $Kq-1$.
For the column bases evaluate at $x=q^{b-1}y$, then
$x=q^{b-2}y$. The same conditions, up to nonzero powers, suffice.
None can vanish in the interior range. Thus these transformations are
defined over $\mathbb Q(q)$ after $K=q^k$ and introduce no denominators
depending on $\theta$ or $X$.

Here are direct verifications of the Fourier completions.
Multiplying the row outside $(x,y)$ factor by $Q^R_{a,j}$ completes
$A(x,q^jy)$. Coefficient division gives $q^{ji}=X^jq^{-j\,\deg_w}$.
The remaining $(z,w)$ roots have indices $b-j,\ldots,k+a-1-j$.
Their common factor is
\[
 H_R=\prod_{u=b-1}^{k-2}v_u,
\]
and the residual degree-$b$ form is \eqref{R22:rowtransform}.
For the column completion, multiplication by $Q^C_{a,j}$ completes
$A(x,q^{b-k-j}y)$. After division and the nonzero character scalar
$X^{k-b}$, it gives $X^{-j}q^{(k-b+j)\deg_w}$.
The outside $(z,w)$ roots become $-a+j,\ldots,k-b-1+j$.
Their common factor and residual form are
\[
 H_C=\prod_{u=0}^{k-b-1}v_u,
 \qquad \mathcal S_b(F).
\]
Thus both transformations represent the actual character quotients
of the cofactor limits, including all shifts in $k$ and $t$.

Let $M_R$ be the $(b+1)\times(b+1)$ coefficient matrix of
\begin{equation}\label{R22:MR}
 z\prod_{u=0}^{b-2}v_u,\quad
 w\prod_{u=0}^{b-2}v_u,\quad
 \mathcal R_b(R_0),\ldots,\mathcal R_b(R_{r-1}).
\end{equation}
Let $M_C$ be the corresponding coefficient matrix of
\begin{equation}\label{R22:MC}
 z\prod_{u=-b}^{-2}(z-Kq^uw),\quad
 w\prod_{u=-b}^{-2}(z-Kq^uw),\quad
 \mathcal S_b(S_0),\ldots,\mathcal S_b(S_{r-1}).
\end{equation}
These are precisely base plus row, and base plus column, after
removing their respective common factors. For the column matrix the
additional nonzero character scalar just used also rescales the base;
its root indices are then $k-b,\ldots,k-2$, as in
\eqref{R22:MC}.

\paragraph{4. Uniform row nonvanishing at every interior bidegree.}
Fix the primitive root $q$ and any interior $k$. We prove
$\det M_R\not\equiv0$ as a polynomial in $\theta,X$ by a specialization
that may use elements of $\mathbb Q(q)$.
In the local bottom-row matrix, retain only nearest forward chords
$c_{i,i+2}=p_i$ and first-column backward chords $c_{h,0}=z$.
Set
\[
 p_i=\ell_1\quad(0\leq i\leq r-3),\qquad
 p_{r-2}=\ell_2,
\]
and set its other chords to zero. Parameters outside this row matrix
can be arbitrary. Its row forms are
\[
 -p_i z\prod_{\substack{1\leq u\leq r\\u\ne i+2}}v_u
 \quad(0\leq i\leq r-2),
 \qquad \ell_r\ell_2z\prod_{u=1}^{r-2}v_u.
\]
For $i\leq r-3$ their normalized forms are
\begin{equation}\label{R22:diagonal_forms}
 -X^2z\prod_{\substack{-1\leq u\leq r-1\\u\ne i}}v_u.
\end{equation}
Evaluate modulo the base space at the $r$ distinct roots
$z/w=q^h$, $0\leq h\leq r-1$. The first $r-2$ columns in
\eqref{R22:diagonal_forms} are diagonal on the first $r-2$ roots and
zero on the last two; their diagonal entries are nonzero for $X\ne0$.

The remaining two local forms are
\[
 U=-\ell_2z\prod_{u=1}^{r-1}v_u,\qquad
 V=\ell_r\ell_2z\prod_{u=1}^{r-2}v_u.
\]
Write
\[
 \alpha=\frac{q-1}{K-1},\qquad
 \beta=\frac{K-q}{K-1},\qquad
 \ell_2=\alpha Q^R_{1,1}+\beta Q^R_{1,2}.
\]
Their normalized forms $P_U=\mathcal R_b(U)$,
$P_V=\mathcal R_b(V)$ are
\begin{align*}
 P_U={}&-z\left[
  \alpha X\prod_{u=0}^{r-2}v_u\,(z-Kq^{-1}w)
 +\beta X^2\prod_{u=-1}^{r-3}v_u\,v_{r-1}\right],\\
 P_V={}&c_1Xz\prod_{u=0}^{r-3}v_u\,(z-Kq^{-1}w)(z-Kw)\\
 &{}+c_2X^2z\prod_{u=-1}^{r-4}v_u\,v_{r-1}(z-Kq^{-1}w)\\
 &{}+c_3X^3z\prod_{u=-2}^{r-5}v_u\,v_{r-2}v_{r-1}.
\end{align*}
where $c_j$ are the coefficients of $\ell_r\ell_2$ in the basis
$Q^R_{2,j}$. Only the first is needed:
\begin{equation}\label{R22:c1}
 c_1=\frac{(1-q^{r-1})(1-q)}{(1-K)(1-Kq)}.
\end{equation}
It follows by evaluating the basis expansion at $x=qy$.

Put $A=q^{r-2}$, $B=q^{r-1}$ and evaluate at $w=1$.
$P_U(A)$ has no $X$ term; $P_U(B)$ and $P_V(B)$ have only their
$X$ terms. The coefficient of $X^2$ in the last two-by-two
evaluation determinant is therefore
\begin{equation}\label{R22:last_minor}
 \begin{split}
 [X^2]\big(P_U(A)P_V(B)-P_V(A)P_U(B)\big)
 ={}&\alpha c_1AB
 \prod_{u=0}^{r-2}(B-q^u)
 \prod_{u=0}^{r-3}(A-q^u)\\
 &\cdot(B-K/q)(A-K/q)(A-K).
 \end{split}
\end{equation}
Every factor is nonzero for the fixed primitive $q$ and every
$2b-1\leq k\leq m-2b+1$. In particular the potential equalities
$K=q^{b-1},q^{b-2},q^{b-3}$ lie outside the range, as do $K=1,q^{-1}$.
This also covers $b=3$, when the first diagonal block is empty.
The kernel of evaluation at these $r$ roots is exactly the
two-dimensional base space in \eqref{R22:MR}. Thus the row matrix
has nonzero determinant as a polynomial in $X$ at this specialization,
for every interior $k$ and every primitive $q$ under consideration.
Its specialized degree in $X$ is at most $2b-3$: the first $b-3$
evaluation columns contribute $X^{2(b-3)}$, and the last determinant
has degree at most three.

\paragraph{5. Column nonvanishing at the same primitive root and bidegree.}
The complete parameter domain in \eqref{R22:types} is invariant under
the reflection of additional positions
\begin{equation}\label{R22:reflection}
 (i,j)\longmapsto(b-j,b-i).
\end{equation}
It preserves $j-i$, hence preserves the forward $(x,y)$ type and
backward $(z,w)$ type. Transpose the global cyclic pencil and reverse
its indices by setting $M'_{uv}=M_{b-v,b-u}$ modulo $m$.
Its diagonal and cycle entries are
\[
 \ell_{b-u},\qquad v_{b-u-1}.
\]
With
\[
 q'=q^{-1},\qquad y'=q^by,\qquad w'=q^{b-1}w,
\]
they become the standard cyclic entries $x-(q')^uy'$ and
$z-(q')^uw'$. This is an invertible change of parameters and variables
over $\mathbb Q(q)$.

More explicitly, if $C$ is the transposed supported-column matrix of
the original block and $B'$ is the supported-row matrix of the new
block, then $B'_{ij}=C_{r-1-i,r-j}$. Reversing both rows and columns
preserves the determinant signs of each square minor. It changes the
row relation $a_iE_i-b_iE_{i+1}$ to the negative of the corresponding
reversed relation. This proves \eqref{R22:Sforms} from
\eqref{R22:Rforms} and identifies the two collections of actual
cofactor limits. An external original column $c=m+g-k$ becomes the
external new row
\[
 b-c\equiv k+b-g\pmod m,
\]
exactly the row alignment in \eqref{R22:alignment}. Bidegree $k$ is
unchanged. The change of $y,w$ only rescales character coordinates by
nonzero numbers; it does not mix them or change their available
coefficient interval.

For precision, the following argument does not assume that separate
generic row and column ranks are already simultaneous. Apply the
row specialization just proved at the primitive root $q^{-1}$ and
at this same $k$. Its determinant is a nonzero polynomial in $X$
of degree at most $2b-3$. There are $T=m-4b+3>2b-3$ interior
characters, so at least one of their distinct root values gives
nonzero determinant. Restricting the resulting binary forms to the
allowed coefficient interval preserves their independence: the
Toeplitz/Schur argument in the next paragraph applies to
$H_R\operatorname{Sym}^b$, using any $b+1$ consecutive allowed
indices. Reflection therefore gives a column specialization at $q,k$
with independent actual column classes. If $M_C$ were identically
singular as a function of its parameters and $X$, such independence
after restriction would be impossible. Hence
\begin{equation}\label{R22:two_nonzero}
 \det M_R(q,q^k,\theta,X)\not\equiv0,
 \qquad
 \det M_C(q,q^k,\theta,X)\not\equiv0
\end{equation}
for every interior $k$, over the same parameter ring and the same
fixed field $\mathbb Q(q)$.

\paragraph{6. Joint row--column rank and coefficient restriction.}
Suppose both matrices are nonsingular for some interior $k,t$.
Before coefficient restriction, base plus row spans
$H_R\operatorname{Sym}^b\mathbb C^2$, and base plus column spans
$H_C\operatorname{Sym}^b\mathbb C^2$. Their common factor is
\[
 H=\prod_{u=b-1}^{k-b-1}v_u,
 \qquad N:=\deg H=k-2b+1\geq0.
\]
After dividing by $H$, the two spaces are
\[
 \left(\prod_{u=k-b}^{k-2}v_u\right)\operatorname{Sym}^b\mathbb C^2,
 \qquad
 \left(\prod_{u=0}^{b-2}v_u\right)\operatorname{Sym}^b\mathbb C^2.
\]
The two degree-$b-1$ factors are coprime, since their indices are
distinct and lie below $m$. Their intersection is exactly their
product times $\operatorname{Sym}^1\mathbb C^2$, namely the base.
Their joint span is therefore all
\begin{equation}\label{R22:joint_space}
 H\operatorname{Sym}^{2b-1}\mathbb C^2,
\end{equation}
of dimension $2b$. This is the required joint independence, not an
addition of two ranks without checking their intersection.

The allowed $w$-degree interval in character $t$ is
\[
 \max\{0,t-(m-k)\}\ \leq j\leq\ \min\{k,t\}.
\]
Its length is at least $2b$, since $k,m-k,t,m-t\geq2b-1$.
Choose $2b$ consecutive indices $e,\ldots,e+2b-1$ in it;
then $0\leq e\leq N$. Restriction of multiplication by $H$ to
these indices has determinant
\begin{equation}\label{R22:schur}
 \det(H_{e+i-j})_{0\leq i,j<2b}
 =\pm s_{((2b)^e)}(q^{b-1},q^b,\ldots,q^{k-b-1})\ne0.
\end{equation}
The identity is dual Jacobi--Trudi, as in R-16/R-19 and R-B15.
For nonvanishing, there are $N$ geometric variables, and the Weyl
numerator exponents for this rectangle are distinct in
$[0,N+2b-1]=[0,k]$. Since $k<m$, their powers of $q$ are distinct,
and the numerator is a nonzero Vandermonde determinant; the
denominator is nonzero for the same reason. For $N=0$ the coefficient
matrix is the identity. This proves that all $2b-2$ normal classes
remain independent modulo the entire base character space.
The same proof with $b+1$ in place of $2b$ and the corresponding
geometric factor proves the row-only restriction fact used in the
preceding paragraph.

\paragraph{7. One integer grid choice for all bidegrees.}
We now fix $m,b,q$ and choose the motif parameters once.
By \eqref{R22:grades} and \eqref{R22:rowtransform}, the $X$ exponents
of the row determinant lie in
\[
 [b-1,2b-1]:
\]
every one of its $b-1$ normal columns has minimum exponent one,
the $b-2$ columns with $a=1$ have maximum two, and the column with
$a=2$ has maximum three. The column determinant has exponents in
$[-b,0]$, by \eqref{R22:coltransform} and the sum
$2+(b-2)=b$ of its $a$ values. Consequently, for each interior $k$,
\begin{equation}\label{R22:Fk}
 F_k(\theta,X)=X\det M_R(q,q^k,\theta,X)
                    \det M_C(q,q^k,\theta,X)
\end{equation}
is an ordinary polynomial in $X$ of degree at most $2b$.
Its coefficients lie in $\mathbb Q(q)[\theta]$, and its total
parameter degree is at most $4(b-1)$. By
\eqref{R22:two_nonzero} it is nonzero. The product over all $T$
interior values of $k$ is therefore a nonzero polynomial with
parameter degree at most $D=4(b-1)T$.

A polynomial of degree at most $D$ in each parameter over a field
of characteristic zero
cannot vanish at all points of $\{0,\ldots,D\}^{2b(b-2)}$ unless
it is zero: prove this successively in each variable by univariate
interpolation. Apply this over the field $\mathbb Q(q)(X)$.
There is an integer grid choice $\theta$ for which the product is
nonzero as a polynomial in $X$. Fix that one choice. Every
$F_k(\theta,X)$ is then nonzero, so at most $2b$ interior characters
are bad for each $k$. At least
\begin{equation}\label{R22:good_modes}
 T-2b=m-6b+3
\end{equation}
characters have all $2b-2$ independent normal classes.

This argument deliberately permits the grid choice to depend on
$m,q$; it makes no claim that one fixed integer motif works for
every root order. It needs neither a bound for $\varphi(m)$ nor a
generic root-of-unity specialization inferred from numerical tests.
If a concrete choice is wanted, enumerate the finite grid in
lexicographic order and test \eqref{R22:Fk} in exact cyclotomic
arithmetic, choosing the first with nonzero product. This is a finite
algorithm, with no polynomial-time assertion.

\paragraph{8. Translations, all bidegrees, and a single differential.}
For fixed $k$, the leading vectors from the block at start $a=bj$
are the transforms by $\tau^{bj}$ of those at zero, where
$\tau:(y,w)\mapsto(qy,qw)$. Suppose a linear combination of these
vectors lies in $W_0$. In a good character $t$, their independent
normal classes give $2b-2$ separate equations of the form
\[
 \sum_{j=0}^{h-1}c_j(q^{bt})^j=0.
\]
Every value of $q^{bt}$ has at most $\gcd(b,m)\leq b$ preimages.
The good characters in \eqref{R22:good_modes} therefore contain
at least $h$ distinct such values. Ordinary Vandermonde independence
forces every $c_j$ to vanish. Hence each of the $T$ bidegrees supplies
$(2b-2)h$ independent leading classes modulo $W_0$.

The actual limits of \eqref{R22:relation} need not be homogeneous
for the $(x,y)/(z,w)$ bigrading. In any relation among them modulo
$W_0$, take the largest selected $(z,w)$ degree. The corresponding
highest components are the independent classes just proved; all
remaining components have smaller degree. Since $W_0$ is bigraded,
these coefficients vanish. Descending in degree proves independence
of all $(2b-2)hT$ full limits modulo $W_0$.

Finally each cofactor combination in \eqref{R22:relation}, multiplied
by its common outside permanent, belongs to the differential image
of the same simultaneous pencil whenever $\varepsilon\ne0$.
Dividing by $2\varepsilon^2$ gives a polynomial family, by
\eqref{R22:local_defect}. Its limit is the full vector just used.
Choose $2m^2+2$ ordinary differential columns independent at
$\varepsilon=0$, and append these polynomial families. The resulting
matrix has a nonzero minor at zero, by the quotient independence.
That minor remains nonzero for a common generic nonzero
$\varepsilon$. The differential rank is therefore at least
$2m^2+2+(2b-2)hT$, proving \eqref{R22:rank}.

\paragraph{9. Quantitative consequences and the precise ceiling.}
Since $h\geq m/b-7$ and $T\geq m-4b$, with both lower bounds
nonnegative under $m\geq8b$,
\[
 J_{m,b}\geq(2b-2)(m/b-7)(m-4b)
 \geq(2-2/b)m^2-22bm.
\]
For $m\geq99$, take $b=\lfloor\sqrt{m/11}\rfloor$.
Then $b\geq3$, $m\geq8b$, and
$\frac12\sqrt{m/11}\leq b\leq\sqrt{m/11}$. Thus
\[
 \frac{2m^2}{b}+22bm\leq6\sqrt{11}\,m^{3/2}
 <20m^{3/2},
\]
which proves \eqref{R22:asymptotic}. The padded image has dimension
$c_m+3$ by R-14, while the determinant image has dimension
$2n^2+2$ by R-17. This gives \eqref{R22:padding} and all the stated
multiplicity consequences. Under \eqref{R22:ratio_threshold},
\[
 4m^2-20m^{3/2}+2\geq2\rho^2m^2+2\geq2n^2+2,
\]
so the dimension comparison is strict after padding. Substitution
$\rho=4/3$ gives the threshold $2025$ and, with G-11/G-12 or the
rectangular specialization of G-17, the claimed exact eventual
rectangular range.

For comparison, the dimension method itself cannot pass the
asymptotic ratio $\sqrt2$. Its necessary strict inequality, using
the upper bound in \eqref{R22:asymptotic}, is
\[
 2n^2<4m^2-2m+3.
\]
In particular it fails for $n\geq\sqrt2\,m$ and $m\geq2$.
The proved lower bound approaches this ceiling from below; for
example $m\geq169$ and
$m<n\leq\sqrt2\,m-4\sqrt m$ suffice by direct squaring.
This does not exclude other methods of finding multiplicity gaps
beyond the dimension comparison, and it does not address
superlinear padding or polynomial-time multiplicity computation.
\end{proof}
\subsection{R-23: the exact homological defect after restriction to four variables}

\paragraph{Statement R-23.}
Let $m\geq3$, let $R=\mathbb C[X_{ij}:1\leq i,j\leq m]$, and let
$J_m$ be the ideal generated by all $(m-1)\times(m-1)$
subpermanents, equivalently the first partial derivatives of
$\operatorname{per}_m$. Let
\[
 S=R/(L_1,\ldots,L_{m^2-4})\simeq\mathbb C[x,y,z,w]
\]
be a general four-variable linear restriction. Write
\[
 \epsilon_m=\dim_{\mathbb C}
       \operatorname{Tor}_1^R(R/J_m,S)_m.
\]
Then
\begin{equation}\label{R23:Tor_defect}
 \dim\mathcal P_{m,4}=4m^2-2m+2-\epsilon_m.
\end{equation}
Moreover
\begin{equation}\label{R23:first_degree}
 \operatorname{Tor}_1^R(R/J_m,S)_d=0\qquad(d<m),
\end{equation}
and R-22 implies
\[
 0\leq\epsilon_m\leq20m^{3/2}-2m\qquad(m\geq99).
\]
Thus attaining the exact row--column torus bound is equivalent to one
specified degree of one Tor group vanishing. This vanishing is not
proved here for $m\geq17$.

The general-section defect is already nonzero in a complete known
range. By R-13, for $3\leq m\leq16$ it is exactly
\begin{equation}\label{R23:small_defect}
 \epsilon_m=4m^2-2m+2-\binom{m+3}{3}>0.
\end{equation}
Consequently the full-variable linear syzygy theorem cannot simply be
carried through a general four-variable restriction. Nor can this
restriction be justified by treating all its defining linear forms as
a regular sequence on $R/J_m$: this is impossible for every $m\geq3$.

\begin{proof}
We first specify the cited full-variable input. Efremenko, Landsberg,
Schenck, and Weyman (2018), Theorem~1.1 and \S4.1, Example~4.1
and Theorem~4.3 in the author version, determine the linear strand
of the subpermanent ideal (R-B16). At subpermanent size $m-1$, the
space of linear first syzygies has dimension $2m-2$. A basis consists
of differences between row Laplace expansions of the full permanent,
and differences between column expansions. In the minimal resolution
of $R/J_m$, therefore,
\begin{equation}\label{R23:resolution_start}
 F_1=R(-(m-1))^{m^2},\qquad
 (F_2)_m=\mathbb C^{2m-2},
\end{equation}
and $F_2$ has no components below degree $m$. Other generators of
$F_2$ have higher degree and do not enter the calculation below.

\paragraph{The old linear relations remain independent.}
Write $G_{ij}=\partial\operatorname{per}_m/\partial X_{ij}$.
The row and column expansion vectors are
\[
 \mathsf R_i=(\delta_{ui}X_{uv})_{u,v},\qquad
 \mathsf C_j=(\delta_{vj}X_{uv})_{u,v}.
\]
Each maps to the permanent under $(f_{ij})\mapsto\sum f_{ij}G_{ij}$.
Thus the $2m-2$ differences
$\mathsf R_i-\mathsf R_1$ and $\mathsf C_j-\mathsf C_1$ are
the cited syzygies. A linear combination is the vector
$((\alpha_i+\beta_j)X_{ij})$, where
$\sum\alpha_i=\sum\beta_j=0$.

After restriction to the cyclic pencil of R-15, both its diagonal
and its cyclic superdiagonal entries are nonzero linear forms.
If this relation vector vanishes, its diagonal entries give
$\alpha_i+\beta_i=0$, and its cyclic entries give
$\alpha_i+\beta_{i+1}=0$. Hence all $\beta_i$ are equal and all
$\alpha_i$ are their negatives. Their zero sums force every
coefficient to vanish. The old relations are therefore independent
on that pencil, and hence on a nonempty Zariski open set of
four-variable restrictions.

\paragraph{There are no new constant relations generically.}
For completeness we prove independence of the $m^2$ restricted
gradient generators on the same cyclic pencil. Number indices
$0,\ldots,m-1$, take a primitive $m$th root $q$, and put
$\ell_i=x-q^iy$, $v_i=z-q^iw$. Each complementary permanent has
one matching. For each $0\leq k\leq m-1$, its $m$ complementary
permanents of $(z,w)$ degree $k$ are the cyclic translates of
\begin{equation}\label{R23:cofactor_orbit}
 F_k=\prod_{i=0}^{m-k-2}\ell_i
                 \prod_{i=m-k-1}^{m-2}v_i.
\end{equation}
For example this is the cofactor deleting row $m-1$ and column
$m-k-1$. Empty products have value one.

Every coefficient of each of the two geometric binary products in
\eqref{R23:cofactor_orbit} is nonzero. Indeed each length is below
$m$, so the geometric binomial-product formula has no factor
$1-q^j$ with $j$ divisible by $m$. For every
$0\leq t\leq m-1$, choose integers
$0\leq a\leq m-k-1$, $0\leq b\leq k$ with $a+b=t$.
The corresponding coefficient of $y^aw^b$ is nonzero. Under
$\tau:(y,w)\mapsto(qy,qw)$ this gives a nonzero component of
character $t$. Thus the orbit
$F_k,\tau F_k,\ldots,\tau^{m-1}F_k$ has dimension $m$ by Fourier
inversion. The different $k$ have distinct bidegrees, proving
independence of all $m^2$ cofactors. Their restrictions are therefore
independent for general $S$.

\paragraph{The exact Tor identity.}
Tensor a minimal free resolution of $R/J_m$ with $S$ and take
homogeneous degree $m$. The part relevant to its first homology is
\begin{equation}\label{R23:degree_complex}
 \mathbb C^{2m-2}\longrightarrow
                 S_1^{m^2}\longrightarrow S_m.
\end{equation}
The first map is injective on the general open set just proved.
The second sends $(f_{ij})$ to
$\sum f_{ij}G_{ij}|_S$. It is exactly the differential of the
permanent projection map, because varying any pencil entry by a
linear form $f_{ij}$ changes the permanent by that expression.
In characteristic zero its rank at a general pencil equals the
dimension $c_m$ of the image closure. Consequently the first
homology of \eqref{R23:degree_complex} has dimension
$4m^2-c_m-(2m-2)$, proving \eqref{R23:Tor_defect}.
In degrees below $m-1$ the module $F_1\otimes S$ is zero. In degree
$m-1$ its map to $S$ is injective by the cofactor independence just
proved. This gives \eqref{R23:first_degree}.
Substitution of R-22 proves the stated upper bound for $\epsilon_m$.

\paragraph{Why the known linear strand does not supply the missing vanishing.}
R-13 gives $c_m=\binom{m+3}{3}$ for $3\leq m\leq16$,
and hence \eqref{R23:small_defect}. The positive values are, in order,
\[
 12,\ 23,\ 36,\ 50,\ 64,\ 77,\ 88,\ 96,\ 100,\ 99,\ 92,\ 78,\ 56,\ 25.
\]
These are general-section Tor groups in the very same permanent
family, even though the complete full-variable linear syzygy space
has only the $2m-2$ displayed generators. The larger-section
vanishing sought for $m\geq17$ needs an additional theorem.
At the special cyclic section the defect is even larger:
\[
 \dim\operatorname{Tor}_1^R(R/J_m,S_{\mathrm{cyc}})_m
 =2m(m-1),
\]
because its differential image has dimension $2m^2+2$ and the old
relations remain independent. This special value is not substituted
for the general-section defect.

Finally Boralevi, Carlini, Micha{\l}ek, and Ventura (2024),
Theorem~4.14, record von zur Gathen's bound
$\operatorname{ht}J_m\geq5$ for every $m\geq3$; their Theorem~4.23
strengthens it to six when $m\geq6$ (R-B17). Therefore
\[
 \operatorname{depth}(R/J_m)
 \leq\dim(R/J_m)\leq m^2-5<m^2-4.
\]
A regular sequence on $R/J_m$ cannot have the required length
$m^2-4$. Thus tensoring its resolution with the four-variable
quotient is not justified as an exact operation by a regular-sequence
argument. This obstruction does not itself decide the single graded
piece in \eqref{R23:Tor_defect}: higher-degree Tor may survive even
if that piece vanishes.
\end{proof}


